\documentclass[fleqn,10pt]{wlscirep}
\usepackage[utf8]{inputenc}
\usepackage[T1]{fontenc}
\usepackage{lineno}
\usepackage{graphicx}
\usepackage{pdfpages}
\usepackage{caption}

\DeclareCaptionType{extendeddatafigure}[Extended Data Figure
][List of Extended Data Figures]
\providecommand*{\theHextendeddatafigure}{\arabic{extendeddatafigure}}
\makeatletter
\newcommand{\saveimageanchor}[1]{%
  \phantomsection
  \expandafter\xdef\csname imagehref@#1\endcsname{\@currentHref}%
}
\newcommand{\labeltoimage}[1]{%
  \begingroup
  \edef\@currentHref{\csname imagehref@#1\endcsname}%
  \label{#1}%
  \endgroup
}
\makeatother

\usepackage{natbib}
%\linenumbers
%\modulolinenumbers[1]
%\title{Hippocampal scaffold of lifelong spatial memory}
\title{\fontsize{18pt}{24pt}\selectfont Persistent hippocampal trace of lifelong spatial memory}
\author[1,*]{Jiang Wu}
\author[1,+,*]{Markus Meister}
\author[1,+,*]{Carlos Lois}
\affil[1]{California Institute of Technology, Division of Biology and Biological Engineering, Pasadena, 91125, USA}

\affil[*]{correspondence to: jiangwu@caltech.edu;meister@caltech.edu;clois@caltech.edu}

\affil[+]{these authors contributed equally to this work}

%\keywords{Keyword1, Keyword2, Keyword3}

\begin{abstract}
\internallinenumbers
Memory presents a biological paradox: the brain changes continuously, yet brief experience can influence behavior for a lifetime. How such memories are stored in brain circuits remains unclear. Here we show that two days in a binary-tree maze leave mice with spatial memories that survive repeated 90-day absences over as much as a year. Across return visits, mice explored the maze more efficiently than when naive and resumed direct runs to reward sooner, both by a factor of 3. Chronic calcium imaging of the same hippocampal CA1 neurons showed that population activity still resembled its earlier pattern, although most neurons were active in different parts of the maze on return. Even a neuron whose spatial tuning in the maze changed could retain its preferred depth in the branching hierarchy of the binary tree. Depth tuning was more stable than tuning to precise position, and the better it survived, the more efficiently mice started to explore after return. Place fields near reward were also preferentially preserved, and more accurate decoding there accompanied earlier sustained reward runs. With repeated experience, spatial representations changed less, both within a visit and across the long absences. The lasting product of repeated experiences was thus not only a trace of the past, but also a preparation for the next encounter.

\end{abstract}
\begin{document}

\flushbottom
\maketitle
\thispagestyle{empty}

\linenumbers
\modulolinenumbers[1]
\section*{Introduction}
Animals often have to turn brief encounters with an environment into long-lasting memories. Such memories must remain useful while the brain that carries them continues to change. A landscape is organized across scales, from broad terrain and paths to junctions and sites that provide shelter or reward \cite{calhoun_ecology_1963, avenant_microenvironment_2003}. To use such a place again, an animal needs to remember where choices leads, which destinations matter, and how to act effectively again without complete recall of the layout. However, commonly used navigation tasks can often be solved using local cues, one or two turns, or a beacon at the goal \cite{wolbers_challenges_2014,jacobs_navigation_2014}. These tasks therefore leave open which features of a complex environment animals remember over long intervals.

The hippocampus is central to spatial navigation in rodents \cite{morris_place_1982}, and it samples convergent input from much of the brain \cite{burwell_cortical_1998}, which makes CA1 an efficient place to ask what an animal has retained. Many hippocampal neurons are active when an animal occupies particular regions of an environment. These regions are called the neuron's place fields, and a neuron's firing pattern across the whole environment forms its place map. These spatial activity patterns have long been considered a candidate neural substrate for spatial memory \cite{okeefe_hippocampus_1971, okeefe_hippocampus_1978}. Yet longitudinal recordings in the same environment show that the place maps of the same hippocampal neurons become less similar over days to weeks \cite{hainmueller_parallel_2018, gonzalez_persistence_2019, lee_statistical_2020, geva_time_2023, climer_hippocampal_2025}, and the synapses onto these neurons turn over on similar timescales \cite{attardo_impermanence_2015,pfeiffer_chronic_2018}. 

Such representational drift and synaptic turnover complicates the understanding of what remains of a spatial memory in the hippocampus \cite{zaki_memory_2025}. Remote memory can become less dependent on hippocampal circuitry \cite{winocur_preserved_2005, squire_memory_2015}, but this does not establish that the hippocampal trace has disappeared. Engram studies distinguish the survival of a trace from its accessibility: hippocampal ensembles can persist in a silent state, no longer recruited by natural retrieval cues \cite{ryan_engram_2015, kitamura_engrams_2017, josselyn_memory_2020, oleary_natural_2024}. Another possibility is that hippocampal stability is selective \cite{sato_distinct_2020, zemla_task-selective_2022, pettit_fos_2022}. Rather than preserving a detailed place field for every location, CA1 may retain the components of a representation that remain most useful for future navigation. 

These accounts focus on what information remains available after the gap. Yet prior experience shapes how a representation forms, biasing which patterns appear \cite{rowland_stable_2011,plitt_experience-dependent_2021,vaidya_formation_2025}. The question is therefore not simply whether hippocampal maps are stable or unstable, but also how maps are formed when the animal returns, and how persistence and rebuilding relate to behavior.

Most studies of representational drift have recorded the same neurons every day or every few days \cite{hainmueller_parallel_2018, gonzalez_persistence_2019, lee_statistical_2020, geva_time_2023, climer_hippocampal_2025}, so repeated exposure to the task may itself maintain or refresh the representations being measured. Behavior in these designs also changes little across sessions, leaving little variation against which representational change can be compared. To study the persistence of the hippocampal code and its relationship with behavior therefore requires three features: a long interval without exposure to the environment, a task that reveals distinct components of spatial memory, and longitudinal recording of the same neurons across the interval.

\section*{Results}

\subsection*{Lifelong memory of a complex navigation task}
We let mice solve a maze task in complete darkness across recurring two-day recording cycles separated by approximately 90~days and spanning much of their adult life (Fig.~\ref{fig:figure1}a,b). The maze was organized as a binary tree embedded in a $15 \times 15$ grid, with six branch decisions separating the stem from each of 64 end nodes \cite{rosenberg_mice_2021}. A fixed water reward was available at one end node and replenished 90~s after each collection. The maze was connected to the home cage, thus allowing animals to enter, explore, drink, and return home without handling. We refer to the first day of initial exposure as the naive day and the first day of each later cycle as a return day. 

On the naive day, mice began by repeatedly revisiting nearby end nodes, then explored broadly for $\sim$ 1~hour before running reliably to the reward (Fig.~\ref{fig:figure1}c, Extended Data Fig.~\ref{fig:ed01}a,b). To distinguish learning from accidental arrivals at the water port, we measured the time to the first long bout, defined as a trajectory to the water port with more than six consecutive correct decisions. Because the home to water path contains only six decisions, a qualifying long bout had to begin at a remote location within the maze, farther from the water port than the home cage. Over the naive day, mice began completing long bouts regularly and continued to do so on the second day (Extended Data Fig.~\ref{fig:ed01}c).

To test whether memory of the task persisted across months, we measured the cumulative time each animal spent in the maze before its first long bout. After the initial 90-day gap, this time fell to roughly one-third of its value on naive day (Fig.~\ref{fig:figure1}d). It decreased further across successive returns at 90-day intervals, up to a year. The number of maze tiles traversed before the first long bout showed the same progression (Extended Data Fig.~\ref{fig:ed01}e-g), indicating that the improvement was not explained by faster locomotion. These results showed a durable behavioral improvement, but did not reveal what the animals remembered. The improved performance upon repeated exposures could reflect knowledge of the maze structure, memory of the reward, or a fixed sequence of turns to reward. 

Across return visits, animals sampled the maze with about three times fewer repeated visits to nearby end nodes than when naive (Fig.~\ref{fig:figure1}c). We quantified exploration efficiency as $32/N_{32}$, where $N_{32}$ is the number of end node visits required to encounter the first 32 unique end nodes. Higher values reflected coverage of half the end nodes with fewer repeated visits (Fig.~\ref{fig:figure1}e, Extended Data Fig.~\ref{fig:ed01}i-k). After sampling a branch, returning mice were more likely to turn back toward the stem, gaining access to other regions of the maze (Extended Data Fig.~\ref{fig:ed02}a-f,i). The increase in stemward turning was greatest near the stem (Extended Data Fig.~\ref{fig:ed02}g). A biased-walk model incorporating this dependence on depth reproduced most of the exploration gain (Extended Data Fig.~\ref{fig:ed02}h).

Mice returning to the maze had experienced both the maze and the reward. To ask what reward experience added, we compared them with latent control animals that had explored the maze without reward one day earlier (Fig.~\ref{fig:figure1}f). When reward was introduced, these control mice began long bouts sooner than naive mice. Returning mice began long bouts even sooner, despite about 90 days away (Fig.~\ref{fig:figure1}g). If recent familiarity with the layout was sufficient to explain the return improvement, latent control animals should have performed at least as well as return animals. We interpret this as evidence of a learned expectation that the maze affords reward, which expressed itself as fewer pauses at decision points before the port was rediscovered and as faster approach once it was (Extended Data Fig.~\ref{fig:ed03}a-d).

Finally, the improvement in getting rewards could still reflect a memorized fixed sequence of left and right choices. We tested this possibility by asking whether mice chose the reward directed branch on their first visit to each of the six decision nodes along the direct path (Fig.~\ref{fig:figure1}h). The assay detected route learning on the second day of a cycle, but after a 90-day gap returning animals were indistinguishable from naive mice, and their small residual bias was explained by an increased tendency to alternate, which latent controls showed equally (Extended Data Fig.~\ref{fig:ed03}e,f). Thus, animals on return days resumed sustained reward runs rapidly, but this behavior did not depend on memory of a fixed motor route.

Taken together, these results demonstrated that after minimal exposure, mice formed maze memories that persisted for almost their entire adult lifespan. This lifelong memory encompassed two distinguishable components: the spatial structure of the maze and existence of a reward.

\subsection*{CA1 keeps a trace after long gaps without maze experience}
We next asked whether this memory had a corresponding neural trace in hippocampal area CA1 over the same long intervals. In transgenic mice expressing GCaMP6s, we used longitudinal one-photon calcium imaging to track dorsal CA1 neurons across recurring recording cycles. The fields of view remained sufficiently stable to match individual neurons across these gaps (Fig.~\ref{fig:figure2}a, Extended Data Fig.~\ref{fig:ed04}a-i), yielding 21{,}412 cross session neuron matches over 36 session pairs from 14 mice, of which 16{,}901 met the activity criterion used in the spatial analyses below.

CA1 population activity at individual maze tiles still resembled the response recorded there 90 days earlier. Across the matched CA1 neurons, the population vector at each location correlated with its earlier counterpart far more than after shuffling neuron identities (mean correlation, 0.22; Fig.~\ref{fig:figure2}h-j). Yet individual place maps were much less faithful. Some retained strikingly similar maps across the gap, whereas others changed almost completely (Fig.~\ref{fig:figure2}b, Extended Data Fig.~\ref{fig:ed05}a). Place map correlations averaged only 0.12, although their mean exceeded the shuffled control and the distribution contained a tail of high correlations absent from the control (Fig.~\ref{fig:figure2}c-e). 

A neuron need not keep its place field to remain one of the more active members of the population. Neurons that were relatively active in one session did tend to remain so 90 days later (mean correlation, 0.53; Fig.~\ref{fig:figure2}m–o). This consistency was stronger in the maze than in the home cage or connecting corridor (Extended Data Fig.~\ref{fig:ed05}e), indicating that the retained signal was specific to task engagement rather than merely reflecting stable intrinsic excitability. Because a population vector registers how active each neuron is, these lasting relative activity rate can carry the pattern at a location across the gap while individual fields move.

Because each recording cycle contained two consecutive days, we next asked whether the long gap persistence depended on which day was compared. All such comparisons spanned approximately 90~days, so any difference between them could not reflect elapsed time. Ensemble rate correlation was highest between the first days of consecutive cycles, indicating that the pattern of relative activity across the population recurred most strongly on re-entry. Population vector correlation followed the same pattern, because this measure also reflects how active each neuron is. Place map correlation, computed within each neuron and therefore insensitive to its overall rate, was instead similar across all day combinations (Extended Data Fig.~\ref{fig:ed05}b-d). 

We next asked for how long the CA1 memory trace could last. We compared the same three stability measures across intervals ranging from 1 day to 3, 6, 9, and 12~months (Fig.~\ref{fig:figure2}f,k,p), with comparisons beyond three months including intervening visits. All three measures were highest across adjacent days and declined sharply across month scale return intervals. Nevertheless, the average correlations remained above the shuffled baseline through the 9-month interval, whereas by 12 months, correlations were no longer distinguishable from chance. Place maps and population measures in CA1 therefore decayed with elapsed time, yet remained detectable across gaps spanning many months.

The wide spread of correlations could arise in two ways. It may reflect a shifting signal, in which different neurons or locations happen to be stable on each return. Alternatively, it may reflect a reproducible organization, in which particular elements of the code are repeatedly more stable than others. The same neurons, maze locations, and relative activity levels that were most stable across one gap tended to remain so across the next gap (Fig.~\ref{fig:figure2}g,l,q). The parts of the representation that changed least therefore tended to be the same from one return to the next.

\subsection*{Experience makes spatial representations more stable}
The preceding analyses compared correlations computed over the entire return days. This was necessary for asking whether any representation survived a months-long gap, but it also compressed the dynamics of the representation that unfolded when an animal re-entered the maze. A whole day correlation could reflect a map that was reinstated immediately at the start of the return day, or it could reflect a representation that was rebuilt during return experience. We therefore measured the time course of map formation upon returning.

An example neuron makes this ambiguity concrete (Fig.~\ref{fig:figure3}a). The neuron recovered a similar spatial pattern after a long absence. However, the place map also changed substantially between the first and second day of each recording cycle. Thus, long gap recurrence and within cycle refinement can coexist in the same neuron.

To measure this refinement directly, we used the second day as a reference and asked how the naive or return day representation approached that reference over time (Fig.~\ref{fig:figure3}b). This moving window analysis followed the same three measures used above: place map, population vector, and ensemble rate correlations. All three measures increased (Fig.~\ref{fig:figure3}e,h,k, Extended Data Fig.~\ref{fig:ed06}c,e,g), indicating progressively closer alignment with the second day reference as animals accumulated time in the maze. Correlations between neighboring windows indicated that the representation changed most rapidly early on the first day and was comparatively stable on the second day, supporting the use of the second day map as the reference (Extended Data Fig.~\ref{fig:ed06}d,f,h).

Yet refinement did not start from scratch. Naive visits were most dissimilar from the second day map, whereas return visits began closer to it and remained more aligned throughout the day (Fig.~\ref{fig:figure3}e,h,k). This improvement was also present when progress was matched by the number of maze tiles visited rather than elapsed time, arguing against a simple sampling explanation (Extended Data Fig.~\ref{fig:ed06}c,e,g). We quantified this improvement across experience by comparing the representation during the first 30~min of each first day with the full second day (Fig.~\ref{fig:figure3}c). Early alignment increased with cumulative experience for all three measures (Fig.~\ref{fig:figure3}f,i,l). Because each cycle produced a different representation, this was not simply recall of one fixed map but a change in where map formation started. 

We next asked whether experience also increased recurrence across recording cycles (Fig.~\ref{fig:figure3}d). Long gap place map and population correlations increased with cumulative recording cycles, whereas ensemble rate correlation changed little beyond the stability already evident (Fig.~\ref{fig:figure3}g,j,m). The neurons' relative activity levels were already stable; what became more consistent with experience was where they were active. The maze itself was unchanged, so a fixed response to the same sensory environment could not by itself explain this increase in similarity.

Cumulative experience also increased with age in this longitudinal design. At the same recording cycles, however, neither 1-day nor 3-month correlations increased with age across mice that first entered the maze between 3 and 10 months (Extended Data Fig.~\ref{fig:ed07}a-f), arguing against age as the basis of the experience effect.

Together, these analyses showed that both the reorganization within visits and the drift between visits diminished with repeated experience.

\subsection*{The reward location anchors selective place stability}
Across successive returns, both the CA1 representation at re-entry and the animals' reward directed behavior improved (Fig.~\ref{fig:figure1}d) in parallel, but these overall trends did not reveal which components of the trace tracked which aspects of performance. The behavioral results revealed two signatures of memory on return: animals resumed long bouts to reward sooner, and they explored the maze more efficiently. We therefore asked whether CA1 carried information specific to either capacity after a long gap.

We trained a spatial decoder on CA1 activity from the first day of one recording cycle and tested it on the first day of the next, 90 days later (Fig.~\ref{fig:figure4}a). Its mean error across the maze was 5.7 maze tiles, only modestly below the 6.6 tiles of the control from shuffling neuron identities (Fig.~\ref{fig:figure4}b). Thus, the preserved code supported decoding of position, but not with precision.

The surviving positional code was not equally reliable throughout the maze. Decoding errors were markedly lower near the reward than at three control locations with matched local geometry (Fig.~\ref{fig:figure4}c, Extended Data Fig.~\ref{fig:ed08}a). Across all session pairs, errors at the control locations were roughly twice those near the reward (Fig.~\ref{fig:figure4}d). This difference held after controlling for spatial sampling and the decoder's spatial prior, and appeared in a population vector correlation analysis as well (Extended Data Fig.~\ref{fig:ed08}b,c). The preserved code was therefore strongest near the reward.

To examine this reward bias along the route itself, we linearized runs between the home cage and the water port and compared, position by position, population activity on Day~$i$ with that on Day~$i+90$ (Fig.~\ref{fig:figure4}e). Because population vector correlation reflects both where neurons fire and how active each neuron is (Fig.~\ref{fig:figure2}), we standardized each neuron's route map to compare spatial preferences relative to its own average and variation along the route (Methods). The population vector correlation peaked at the water port in both directions, exceeding the mean of the other route positions by 0.14 on runs toward the port and by 0.11 on runs away from it (Fig.~\ref{fig:figure4}e). More neurons had fields near the port on runs toward it (on average 18\% versus 12\% of each animal's neurons in the last two route bins, spanning 7.6~cm), but this concentration alone did not produce the population correlation: neuron ID shuffles preserved the concentration and gave correlations near zero throughout. The advantage did not differ significantly between the two directions (runs away minus runs toward, $-0.05$; 95\% CI, $-0.12$ to 0.02). The reward end advantage was positive in 13 of 13 animals on runs toward the port and 11 of 14 on runs away from it. It was also present on the second day of each visit, in recordings not used for Fig.~\ref{fig:figure4}e (0.11 and 0.09; Extended Data Fig.~\ref{fig:ed08}d). It remained after accounting for within-day reliability and after equalizing the number of sampled frames across route positions, and held on routes to and from two other junctions (Extended Data Fig.~\ref{fig:ed08}e-h). On runs toward the port, it was present on both rewarded and timeout visits (Extended Data Fig.~\ref{fig:ed08}i). At the port, approach patterns on one day resembled departure patterns 90 days later, and vice versa (0.17 versus 0.02 at other positions; Extended Data Fig.~\ref{fig:ed08}j), so a preserved component was expressed during opposite directions of travel. Preferential preservation near reward was therefore observed during both approach and departure.

Finally, we asked whether the strength of the preserved code related to behavior. Decoding error across the whole maze predicted how soon mice resumed long bouts, as expected for a behavior that requires traversing a long path (Extended Data Fig.~\ref{fig:ed10}a). The relationship held even at the single tile beside the reward, where lower error accompanied an earlier first long bout (Fig. \ref{fig:figure4}f).

Mice explored about three times as efficiently on return as when naive, yet their earlier place maps supported only imprecise position decoding (Fig. \ref{fig:figure1}e and Fig. \ref{fig:figure4}b). What spatial knowledge could support so large an improvement?

\subsection*{A compact structural code predicts memory guided exploration}
Efficient exploration may not require a precise coordinate code for every maze location. In the binary-tree maze, a more compact variable is tree level, the animal's depth in the branching hierarchy (Fig.~\ref{fig:figure5}a). Level is relevant to exploration because the cost of a wrong turn depends on it: re-entering a sampled branch near the leaves costs only a few extra end node visits, whereas doing so near the stem sends the animal back through a large region it has already covered. 

For each neuron, we computed mean activity at each of the seven tree depths regardless of where in the maze they lie (Fig.~\ref{fig:figure5}a,b); we call this tuning curve the neuron's level map. We correlated these level maps across recording cycles to ask whether neurons preserved their tuning to level after a long gap.
 
As described before, place map correlations were low on average, with most neurons near zero and a long tail of high values (Fig.~\ref{fig:figure2}d). Level map correlations instead peaked near one, so preservation was the typical outcome rather than the exception (Fig.~\ref{fig:figure5}c). The mean level map correlation was significantly higher than a shuffle control (Fig.~\ref{fig:figure5}d), and neurons whose preferred level did shift usually moved to adjacent levels rather than remapping randomly (Extended Data Fig.~\ref{fig:ed09}a,b). In the same neurons, depth tuning was more stable than tuning to precise location.

Because level maps are coarser than place maps, this stability could reflect averaging the maze into fewer bins. We therefore compared three partitions matched at the same number of bins: one following tree level, and two obtained by combining maze tiles in the same subtree as a bin. The level partition was more stable than either alternative (Extended Data Fig.~\ref{fig:ed09}d-f). 

Level and place coding overlapped without being equivalent. Shuffling spatial bins within each level, which destroys detail below the level while leaving level structure intact, halved the place map correlation (Extended Data Fig.~\ref{fig:ed09}c). Roughly half of the place map similarity across the gap was carried by level structure, and half by detail within levels. The two codes nevertheless came apart across neurons. A neuron that preserved its precise field necessarily preserved its level, whereas the converse need not hold, and it was the asymmetric case that dominated: many neurons kept their preferred level while losing place map correlation, and few showed the reverse (Extended Data Fig.~\ref{fig:ed09}g).

This structural memory was also evident at the population level. A level decoder trained before the 90-day gap recovered the animal’s tree level after the gap well above chance (Fig.~\ref{fig:figure5}e). Removing neurons with stable place maps reduced level decoding but did not abolish it (Extended Data Fig.~\ref{fig:ed09}h), so neurons whose place maps reorganized still carried level information. 

What survived in CA1 strongly predicted how well the animal navigated months later: the better the level code was preserved, the more efficiently it explored on return (r=0.75; Fig.~\ref{fig:figure5}f). This single neural measure accounted for more than half of the variation in exploration efficiency (Extended Data Fig.~\ref{fig:ed10}b). Two comparisons showed that this link was specific. Ensemble rate correlation, the most persistent of our measures, did not predict exploration efficiency (Extended Data Fig.~\ref{fig:ed10}a). And among the turning policies that make up exploration efficiency, level map correlation was associated with one alone: the tendency to turn back towards the stem after finishing a branch, which made the largest contribution to the improvement in coverage (Extended Data Fig.~\ref{fig:ed02}d, Extended Data Fig.~\ref{fig:ed10}c). The value of turning back toward the stem depends on the scale of the maze opened up by that choice. A preserved level code carries information about that scale. 

Across three months of absence, detailed place maps reorganized, yet CA1 retained the two variables that made the maze navigable: where reward lay, and how much of the maze a choice opened up.

\section*{Discussion}
Two days in a maze gave mice a memory that outlasted repeated 90-day absences across much of their adult lives. On each return they searched more efficiently than when naive and resumed long bouts to the water port sooner. In CA1, what survived was selective: most detailed place maps had changed, while a compact code for depth in the branching hierarchy and stable coding near the reward persisted, and both were naturally expressed during navigation and related to how the animals behaved. With repeated experience, spatial representations changed less during a visit and across the months between visits.

A memorized sequence of turns would solve this maze. By the second day of every cycle, the mice’s first choice at most of the junctions on the direct path already favored the reward branch (Extended Data Fig.~\ref{fig:ed03}f). After a 90-day gap, that bias was gone. The assay samples the first encounter with each junction after return, so route knowledge may have survived in a form that re-exposure was needed to express. Either way, the turn sequence was not what the memory led with. Absences of this length are ordinary for a wild animal, and over such intervals the world does not hold still \cite{dingle_migration_2014, geiser_seasonal_2020}. A fixed route can fail when a corridor is blocked or the reward moves. What returned across each gap was not the answer but the means of finding it again.

That distinction points to what a durable memory of a place would usefully contain. Natural environments confront an animal with the same kinds of situation over and over in different places: a junction where a path divides, a stretch of open ground, a narrowing that leads to shelter. What defines each kind of situation is what it offers rather than where it is. Knowledge of that kind is compact and carries from one instance to the next \cite{singer_experience-dependent_2010, behrens_what_2018, whittington_tolman-eichenbaum_2020}. If a memory of a place decays unevenly, this is the part one would expect to last \cite{richards_persistence_2017}.

Depth in the branching hierarchy is a variable of this kind, a state abstraction over the many locations that share it \cite{mckenzie_hippocampal_2014}. The preserved tuning could reflect depth itself or one of the features that covary with it in our maze: corridor scale, the number of branches ahead, or distance from the stem. What matters is that such information can guide a choice without a specific memory of the junction: knowing how deep a junction sits is enough to know roughly what a turn will cost, even at a junction not visited for months. Coordinates do not generalize this way. Spatial precision is expensive, and worth keeping only for the places that earn it, such as the reward \cite{sato_distinct_2020, grosmark_reactivation_2021, pettit_fos_2022}. The pattern of relative activity at the port was preferentially preserved whether the animal was arriving or leaving, so the stable code was not tied to the act of approaching the destination. Such a local code could help the animal recognize the destination when encountered without specifying the complete sequence of turns needed to reach it.

Depth and reward are not two memory systems but two scales of one representation. Their associations with behavior differed in strength, but did not separate cleanly. Level map correlation was the better predictor of efficient exploration, and most of what it explained was its own rather than shared with reward code. Reward decoding led for the return of long bouts comparing to level map, but not decisively, since reaching the water port requires knowing the maze as well as the goal (Extended Data Fig.~\ref{fig:ed10}a,b). 

These codes recurred across months, well beyond the timescales of CA1 spine turnover reported in previous studies \cite{attardo_impermanence_2015, pfeiffer_chronic_2018}. The question is then not whether CA1 drifts but what stays fixed while it does. If a representation is worth holding only where it counts, change should collect where it does not, much as the arm can move in different ways and still bring the hand to the same target  \cite{scholz_identifying_2000}. On a linear track, changes in detailed place fields need not prevent an animal from finding a reward at the end. Their instability may therefore reveal how little spatial detail the task demands, rather than how little the hippocampus can preserve \cite{kentros_increased_2004, chen_singlefield_2026}.

Earlier experience shaped not only what persisted across an absence, but also how much the representation need to change when the animal returned. Past learning may restrict the patterns of activity a familiar task can evoke, as proposed in accounts of preconfigured network dynamics \cite{mckenzie_preexisting_2021, buzsaki_neurophysiology_2022} and schemas \cite{tse_schemas_2007}.

The larger opportunity is what can be measured across a long absence. Remote memory has been studied for decades in relatively simple conditions \cite{frankland_organization_2005,squire_memory_2015}, and tagged c-Fos ensembles have shown that hippocampal neurons can be reactivated \cite{tayler_reactivation_2013, tanaka_hippocampal_2018, refaeli_engram_2023, ko_systems_2025}. What an overlap between populations cannot say is what the trace contains. A representation can preserve some features while changing others and following the same neurons and their tuning across months makes that selectivity visible. Applied to other circuits, this approach could reveal whether information that fades from CA1 persists elsewhere. The next challenge is to understand how a changing brain preserves what an animal has learned and revises it when a familiar place presents a new problem.

\newpage
\section*{Methods}

\subsection*{Experimental procedures}

\subsubsection*{Animals}

All procedures were approved by Institutional Animal Care and Use Committees at Caltech under protocols 1656 and 1709. Experiments used C57BL6J-Tg-Thy1-GCaMP6s mice (Jackson Labs stock: 025776) \cite{dana_thy1-gcamp6_2014}. A total of 36 mice (17 female and 19 male) were used in this study. Of these, 18 contributed at least one return session after a gap of approximately 90 d, and 22 contributed dorsal CA1 calcium imaging, of which 14 had cross session registration maintained across long gaps. See \hyperlink{supp-table-1}{Supplementary Table 1} for a breakdown of which mice were used in which experiments and analyses. Mice were 1 to 7 months old at surgery and 3 to 10 months old at their first maze exposure. They were single housed post surgery on a 12 hr/ 12 hr dark/light cycle and experiments were performed in the dark phase. When mice received water restriction, their body weights were monitored every day. During the intervals between recording cycles, mice were housed in a different room and had access to water and food ad libitum. 

\subsubsection*{Behavior task}

The maze formed a binary tree with six branch choices between the home entrance and each of 64 end nodes. It was connected to the home cage, allowing mice to enter and leave freely throughout a session. A water port was installed at one end node, which was connected to the entrance by a unique direct route. The rewarded end node was fixed and remained the same across recording cycles. Maze construction followed the previously described design \cite{rosenberg_mice_2021}. The apparatus had an approximately $24\times24$-inch footprint and was laid out on a 1.5-inch square grid. Walls were made from 1/8-inch infrared transmitting acrylic (ePlastics) that was opaque in the visible spectrum, and the floor was infrared transmitting acrylic supported by a thicker clear acrylic sheet. The resulting corridors were approximately 1-1/4 inches wide. The walls were 5 inches high and the maze had no ceiling, leaving the top open to accommodate the miniscope tether. Experiments were performed in total darkness under diffusive 940 nm IR illumination (CMVision) for behavior recording. The floor and walls were first wiped with disinfectant wipes (Rescue) and then washed with ethanol and water after each session. A nose poke at the water port triggered delivery of 20 $\mu L$ water through a steel tube controlled by a Bpod state machine (Sanworks LLC.). After each delivery, the port remained inactive during a 90 s timeout. Recording lasted for 6 hrs for naive days and 4 hrs for all the other days. Reward also triggered an infrared LED recorded by the behavioral camera to align with the behavior video.

Each recording cycle comprised two sessions separated by 24~h. Day~$i$ of the first cycle was the animal's naive session; subsequent Day~$i$ sessions were return sessions conducted after approximately 90~d without maze experience (90~days $\pm$ 4~days per interval with range 82 to 110 d) with actual interval lengths summarized in \hyperlink{supp-table-2}{Supplementary Table 2}. The corresponding Day~$i+1$ session, termed the second day session, was conducted 24~h later. Water source (hydrogel) was removed from the home cage 20~hrs prior to each recording cycle and food pellets were provided on the home cage floor with bedding materials.  

In the latent condition, the water port was replaced by a maze wall on latent day1, so mice explored the complete labyrinth without a reward port. The wall was removed and a water port was installed later before the start of latent day2, when water became available. The two sessions lasted 6 hrs and 4 hrs, respectively. On latent day1, food pellets and hydrogel were provided in the home cage. Between sessions and on latent day2, mice were water restricted and received only food pellets.

\subsubsection*{Surgery}

All surgical procedures were performed in aseptic conditions. Animals were anesthetized via isoflurane induction chamber at 5\% concentration and then placed in a stereotaxic frame on isoflurane at a concentration between 0.5\%-2.5\% with body temperature maintained at 37~$^\circ$C. Ketoprofen (5~mg~kg$^{-1}$) and buprenorphine (3.25~mg~kg$^{-1}$) were administered subcutaneously before surgery. Bupivacaine (1~mg~kg$^{-1}$) was applied along the incision before closure, and ibuprofen was provided in the drinking water for at least 3~d after surgery at 30~mg~kg$^{-1}$.

GRIN lens implantation directly dorsal to CA1 followed the procedure described previously \cite{gonzalez_persistence_2019}. A 1.0~mm or 1.8~mm diameter craniotomy was centered 2.0~mm posterior and 1.8~mm lateral to bregma on the right hemisphere. Fresh artificial cerebrospinal fluid was applied throughout the procedure. Using a blunt 26-gauge needle under continuous irrigation, the dura, overlying cortex and part of the corpus callosum were aspirated until a thin layer of horizontal fibers became visible and remained intact over the hippocampus. A GRIN lens (1.8~mm diameter, 0.25 pitch, 0.55 NA, Edmund Optics \#64-519 for the custom miniscope and 1.0~mm diameter, 0.5 pitch, 0.5 NA, Inscopix Inc. 1050-002176 for the UCLA miniscope V4) was lowered to 1.35~mm below the skull surface. The lens and two anterior skull screws were secured with cyanoacrylate and acrylic dental cement, and the exposed lens was protected with Kwik-cast.

After at least a month of recovery, mice were anesthetized with isoflurane and a miniscope was aligned to the GRIN lens under live fluorescence to select the focal plane. For animals recorded with the custom miniature microscope, the device was permanently fixed to the implant with dental acrylic and its focusing mechanism was sealed with cyanoacrylate. For animals recorded with the UCLA miniscope, a baseplate was fixed to the implant with dental acrylic and the UCLA miniscope was screwed to the baseplate one day before every recording cycle. The focal plane was relocated each time the UCLA miniscope was reattached using the movie recorded during the previous recording cycle as template. Mice recovered for at least a week after custom miniscope or baseplate attachment before the first behavior experiments.

\subsubsection*{Calcium and behavioral recording}

For 16 animals, hippocampal calcium activity was recorded with a custom miniscope constructed as described previously \cite{gonzalez_persistence_2019}. The microscope used the sensor board from an MC900 miniature camera containing an OmniVision 7960 CMOS sensor; the sensor board together with the clock and voltage regulator board were connected to a USB analog video to digital converter. LED intensity was controlled through the DAC output of a signal generator. The optical assembly was housed in a black acetal body (Delrin) made by a 5-axis CNC machine (PocketNC) and contained a 470/40~nm excitation filter, a 495~nm long-pass dichroic mirror, a 520/50~nm emission filter (Chroma technologies) and a 12.5~mm achromatic lens (Edmund Optics). The field of view was approximately $600\times479~\mu$m at $720\times576$ pixels. Another 6 animals were recorded with UCLA miniscope V4 \cite{aharoni_lab_and_ucla_miniscope_project_miniscope_nodate}.

For custom miniscope, calcium movies were acquired at 25~Hz using USB 3.0 video capture device (StarTech.com, USB3HDCAP). For UCLA miniscope V4, calcium movies were acquired at 25~Hz using Miniscope DAQ. The movies were saved in lossless FFV1 codec format. Excitation power was between 0.05 to 0.2 mW and fixed for each animal and the excitation LEDs were turned off most of the time when animals were outside of the maze. A customized miniature commutator system was added between the miniscope and the computer and attached to a pulley system so that the animals could easily turn and travel with the tether. 

Behavior was filmed with Grasshopper3 (GS3-U3-23S6M, Flir) and V0828-MPY2 Lens (8~mm, 1.1", Computar) at 25~Hz and $1,920 \times 1,200$ pixels. Behavioral and calcium recordings were synchronized at 25~Hz using Bonsai \cite{lopes_bonsai_2015}, with a maximum error of 20 ms. 


%\subsubsection*{Histology}

%After terminal anesthesia with Katamine, mice were transcardially perfused with 4\% paraformaldehyde in phosphate-buffered saline (PBS; pH~7.4). Brains were removed, post-fixed overnight at 4~$^\circ$C in the same fixative, embedded in 3\% agarose and sectioned coronally at 80~$\mu$m on a vibratome. Sections were washed in PBS for 30~min, permeabilized for 15~min in PBS containing 0.3\% Triton X-100, washed again and blocked for at least 1~h at room temperature in PBS containing 10\% fetal bovine serum. Sections were incubated overnight at 4~$^\circ$C with chicken polyclonal anti-GFP antibody (1:1,000; Aves Labs, GFP-10-10) and mouse monoclonal anti-RGS14 antibody (1:250; UC Davis, 75-170) in PBS containing 1\% fetal bovine serum. After three PBS washes, sections were incubated for 90~min with Alexa Fluor 488 goat anti-chicken and Alexa Fluor 555 anti-mouse secondary antibodies (1:1,000; Invitrogen) in PBS containing 1\% fetal bovine serum and mounted with Fluoromount. Sections were imaged with xxx microscope. 

\subsection*{Data preprocessing}

\subsubsection*{Position tracking}

DeepLabCut was used to estimate seven body keypoints (nose, tail, center of mass, and four paws) in each frame using a network trained on approximately 3,000 hand labelled frames \cite{mathis_deeplabcut_2018}. For each recording, four manually annotated maze corners were used to estimate a projective transformation from image coordinates to normalized maze coordinates. Positions were discretized on the $15 \times 15$ maze grid. Coordinates with likelihood below 0.95 were removed and linearly interpolated from the remaining detections. Transitions incompatible with the maze topology were repaired between valid anchors using shortest wall respecting paths. Decision node visits were registered upon nose entry into the designated tile of a junction or end node; re-entry into the same node before visiting another node was not counted as a new visit. Nose position was used for route progress and junction choices because it provided the earliest unambiguous entry into a selected branch. Center of mass was used for speed, spatial activity maps and alignment with calcium activity. 

Unless specified otherwise, position dependent neural analyses used speeds greater than 2 and less than 60~cm~s$^{-1}$. For behavioral speed estimates, continuous center of mass coordinates were median smoothed over five frames within each uninterrupted trajectory. Segments faster than 60~cm~s$^{-1}$ were discarded as tracking artifacts. For behavior mode specific speed, speed was the retained distance divided by the time assigned to that state. 

\subsubsection*{Calcium imaging preprocessing}

Raw calcium videos were transcoded offline to H.264/MPEG-4 AVC in an MP4 container using lossless settings (libx264, preset medium, qp 0, yuv420p). Videos were then temporally downsampled 4× by averaging non-overlapping blocks of 4 consecutive frames, implemented as two successive stages of pairwise frame averaging and decimation. Two recording sessions from the same recording cycle were concatenated and final analysis files were encoded at CRF 10. All procedures were performed in ffmpeg through custom MATLAB (2023b) scripts. Downsampled calcium movies were motion corrected with NoRMCorre \cite{pnevmatikakis_normcorre_2017} using rigid registration within a manually selected rectangular mask. Shifts were estimated from a spatially band-pass filtered copy of the movie using an expected cell radius of 10 pixels. Motion corrected movies were spatially downsampled 2x for calcium signal extraction. 

Spatial components and fluorescence traces were extracted with EXTRACT \cite{dinc_fast_2024}. Extraction used an expected cell radius of 4 pixels, spatial high- and low-pass cutoffs of 6 and 1. Components were screened using a trace signal-to-noise ratio greater than 7, a temporal-corruption score below 0.7, a footprint area of 0.1--6 times the expected area, a spatial-corruption score below 5, and an eccentricity below 8. Candidates were subsequently semi-manually curated with ActSort \cite{jiang_actsort_2024}. Human annotations took precedence for $\sim100$ reviewed candidates for each movie, whereas classifier predictions supplied labels for the remaining unreviewed candidates. Fluorescence traces were deconvolved with the OASIS algorithm \cite{friedrich_fast_2017}. Deconvolution used FOOPSI with a first-order autoregressive, AR(1), calcium-dynamics model. Noise was estimated separately for each trace from its high-frequency power spectrum. The decay coefficient was estimated from the trace autocovariance, and the decay and baseline parameters were optimized. The minimum event amplitude was five times the estimated noise standard deviation ($smin = -5$), with no additional $\ell_{1}$ penalty ($\lambda = 0$). Deconvolved activity was converted back to the 25 Hz acquisition rate by repeating each downsampled value across its four constituent frames for alignment with behavior, so all frame counts and integration windows refer to 25 Hz samples.

\subsubsection*{Longitudinal cell registration}

Day~$i$ and Day~$i+1$ movies processed in a single fit shared component identities by construction. For all other comparisons, spatial footprint projections were first aligned by translation and then refined by nonrigid demons warping \cite{thirion_image_1998}. Primary longitudinal identities were constructed from pairwise CellReg registrations \cite{sheintuch_tracking_2017} between consecutive recording cycles that were generally separated by approximately 3 months. These intervals typically retained sufficient shared field of view. Pairs for which the two-stage alignment left insufficient overlap for reliable cell matching were excluded from cross-session analyses. Candidate pairs within 14~$\mu$m were evaluated using CellReg's spatial correlation model. Final assignments were made by probabilistic clustering at $P_{\mathrm{same}} \geq 0.5$. Identities spanning 6-, 9-, or 12-month intervals were obtained by transitively composing adjacent cycle maps through their shared intermediate session identities. An identity was retained only when a complete one-to-one chain could be followed through every intervening session. 

\subsection*{Behavioral analysis}

\subsubsection*{Behavior mode}

Behavioral modes were assigned from the reconstructed node trajectories. Behavior in the maze was divided into three mutually exclusive modes: "drink", "leave" and exploration. The "drink" mode was defined as direct routes to the reward with all correct turns in rewarded sessions. The "leave" mode was direct routes to the home with all correct turns. All remaining intervals in the maze were classified as exploration. 

\subsubsection*{Reward acquisition and locomotor control}

We quantified acquisition by when animals first encountered reward and when they began sustained reward directed navigation. Reward acquisition was described by cumulative reward count. Sustained reward directed navigation was measured by cumulative long bouts. A long bout was an uninterrupted sequence of more than six consecutive transitions between adjacent locations in the maze topology, each reducing shortest path distance to the water port. First long bout time is measured at the onset of the qualifying path, not its arrival at the port. Progress to the first reward, tenth reward and first long bout was measured as cumulative time in maze and maze tile transitions. Cumulative reward and long bout counts were evaluated on 10~s time grids or 50 transition grids. The two progress measures were transformed as $\log(1+x)$ for comparisons, whereas regressions against neural measures used elapsed seconds.

\subsubsection*{Directed approaches}

Reward approaches comprised all long bouts to rewarded end node that began after the first reward collection. Control approaches followed the same "more than six correct turns" rule but ended at the other three unrewarded end nodes at symmetric locations. Eligible control approaches were pooled within session. Reward approach speed advantage was defined as
\begin{center}
\(
\log\left(\frac{v_{\mathrm{reward}}}{v_{\mathrm{control}}}\right),
\)
\end{center}
where $v_{\mathrm{reward}}$ and $v_{\mathrm{control}}$ were the median reward- and pooled control-approach speeds, respectively.

\subsubsection*{Exploration efficiency}

We scored exploration by how much redundant end node sampling an animal required while covering the maze. An arrival at an end node counted as a visit, and a new visit was counted only after the animal left and returned. If $N_{32}$ was the visit on which the 32nd unique end node was reached, exploration efficiency was
\begin{center}
\(
E_{32}=\frac{32}{N_{32}}.
\)
\end{center}

Repeated coverage was measured in coverage rounds ending when 32 distinct end nodes had been visited; counting restarted at the next end node entry, and incomplete final rounds were omitted. Time resolved round efficiency was summarized in 10~min or 1,000~transition windows. A session contributed one equally weighted value when it completed a round in that window, and a window required at least five contributing sessions. 

\subsubsection*{Junction turning policies and biased walk simulations}

We used local turning policies to ask which recurrent choice tendencies were sufficient to reproduce the reduced sampling redundancy of return sessions. Policies were calculated from transitions at non-root internal junctions during exploration mode. After arrival from the stem, $P_{\mathrm{SF}}$ was the probability of entering either branch instead of reversing. Conditional on continuing, $P_{\mathrm{SA}}$ was the probability that the outgoing branch had the opposite left--right orientation to the preceding stem. After arrival from a branch, $P_{\mathrm{BF}}$ was the probability of leaving that branch instead of immediately returning to it, and $P_{\mathrm{BS}}$ was the probability of continuing toward the stem rather than entering a branch. 

Turning policies were estimated during the first 30~min of cumulative time in maze. Junction weighted estimates calculated each probability separately at every eligible junction and then averaged the junction values. Junction weighted estimates were used for the four-policy comparisons and corresponding biased walk simulations. For depth-varying $P_{\mathrm{BS}}$, decisions were pooled within three depth groups 1--2 (shallow), 3--4 (middle) and 5 (deep). The combined depth-weighted estimate assigned a depth-$d$ decision weight $2^{6-d}$, proportional to the number of end nodes below that junction. Analyses comparing all four junction policies included sessions with at least 30 decisions and at least 10 junctions contributing to each policy estimate. Depth-varying $P_{\mathrm{BS}}$ analyses instead required at least 30 decisions and at least 10 decisions in each depth group.

We simulated walks on the whole binary tree using the measured policies. The depth-varying model used separate shallow, middle and deep $P_{\mathrm{BS}}$ values. At the root decision, the two branches were chosen with equal probability. At subsequent junctions, the current and immediately preceding locations identified the incoming edge as a stem or branch. Transitions then followed the corresponding junction policy, and each end node visit was followed by a return to the preceding junction. Each simulated sequence ended after the number of end node visits observed in the corresponding session and was scored with $E_{32}$. One hundred walks were averaged for each session. 

We asked how much of the observed improvement in exploration efficiency from naive to return sessions could be reproduced by the measured changes in junction policies. For each return session, we simulated exploration using either its own turning policies or those measured when the same animal was naive. The mean simulated difference was divided by the observed difference between return and naive sessions. To estimate the contribution of each policy, one return day turning policy was replaced by its value measured when the same animal was naive while the other policies were held fixed. The contribution was calculated as the simulated efficiency before replacement minus the efficiency after replacement.

\subsubsection*{Junction dwell and first route choices}

Junction dwell analyses tested whether prior experience changed the time animals spent at the junctions before and after encountering reward. Dwell was measured for each completed choice at a junction. Each junction was classified as lying on or off the direct home to reward route.

The measurement window began when the animal arrived at the junction and ended when it reached the next junction. Within this window, dwell was the total time that the animal's body center occupied only the junction tile. Time spent traveling along the connecting branch was therefore excluded. If the animal leaves and returns to the junction before reaching the next junction, that later junction occupancy is also counted. Therefore, the metric is total junction occupancy during a completed choice, not necessarily one uninterrupted dwell. Measurements with dwell greater than 0~s and no more than 30~s were retained.

First visit route score was the fraction of the six decision junctions on the direct reward route at which the animal's first choice directed towards reward. Each session therefore contributed one route score, ranging from 0/6 to 6/6. The first choice at each junction was also examined separately as a binary reward directed or non-reward directed choice. A session contributed one such value at each of the six junctions.

\subsection*{Neural analysis}

\subsubsection*{Neural activity maps and stability measures}

We quantified longitudinal persistence at the levels of individual spatial maps (place maps), location specific population patterns (population vectors), ensemble recruitment (ensemble rates) and tree depth (level maps). Unless stated otherwise, registered neurons were included when their deconvolved event rate was at least 0.1~min$^{-1}$ in both sessions.

For each retained sample $t$, $a_t$ denoted denoised activity and $\Delta t$ denoted the duration represented by that sample. For place maps, $x_t$ identified a spatial bin on the $15\times15$ maze grid. For level maps, $x_t$ identified one of the seven tree depths (0--6). The occupancy normalized activity at spatial bin or depth $j$ was
\begin{center}
\(
\lambda_j =
\frac{\sum_{t:x_t=j} a_t}
{\sum_{t:x_t=j} \Delta t}.
\)  
\end{center}
Thus, $\lambda_j$ represented activity per unit occupancy time in spatial bin $j$ for a place map or at depth $j$ for a level map. All quantitative spatial analyses used the $15\times15$ maze grid for robust estimation. Selected place maps were rendered on a $45\times45$ grid only to visualize finer spatial structure; these higher resolution maps were not used for inferential analyses.

Place maps contained one neuron's activity across traversable maze locations, and level maps contained its activity across tree depths 0--6. Spatial bins or depths were included when occupied for at least 5~s in both sessions. Maps were not spatially smoothed; in particular, no two-dimensional smoothing was applied that could mix locations separated by maze walls. Place map and level map correlations were calculated for each registered neuron and then averaged within the session pair. A population vector comprised the activity of all registered neurons at one spatial bin. Population vector correlations were calculated across matched neurons and then averaged across shared spatial bins. An ensemble rate vector contained one value per registered neuron, calculated as total denoised calcium activity divided by the analyzed duration. Ensemble rate correlations were calculated between vectors containing all matched neurons, including neurons that were silent in one of the two sessions. All correlation analyses used Pearson correlation.

For correlation nulls, later-session neuron IDs were permuted independently within each session pair in 10,000 repetitions while maps, occupancy and population size were held fixed.

\subsubsection*{Long gap correlation and recurrence}

We examined how neural correlation varied with the time between sessions by grouping session pair values according to interval. Only Day~$i$ was considered from each recording cycle in this analysis.

Recurrence was assessed in three-session sequences by relating correlation across the later gap to correlation across the preceding gap. Statistical models used the underlying unbinned observations rather than the binned values shown in the figures.

Long gap pairs were also grouped according to whether the source and target were Day~$i$ or Day~$i+1$ sessions. Day~$i$-to-Day~$i$ correlation was compared with the other three combinations after averaging session pair values within animal.

\subsubsection*{Ensemble rate correlation during and outside maze behavior}

We compared ensemble rate correlation during maze behavior with correlation outside the maze using Day~$i$ to Day~$i+90$ pairs. Outside maze samples combined all frames recorded in the home cage or connecting tunnel. Within each session pair, same number of frames as outside maze samples were drawn from samples inside the maze without speed filtering. Ensemble rate correlations were averaged across 20 draws.

\subsubsection*{Within cycle convergence analysis}

We used the full Day~$i+1$ representation as a fixed reference to measure how the Day~$i$ representation developed during renewed maze experience. Place maps, population vectors and ensemble rates were calculated in overlapping 10, 20 or 30 min Day~$i$ windows advanced by 60~s and in 3,000 transition windows advanced by 100 maze tile transitions. Level maps were calculated in overlapping 30 min windows advanced by 60~s. Each representation was correlated with the corresponding full Day~$i+1$ representation. Each window contained at least 300~s of calcium data. 

For each window definition, $n$ was the number of nonoverlapping windows of that span that fit within the session, calculated from cumulative time in maze or cumulative maze tile transitions as appropriate. A spatial location required at least $5/n$~s in the Day~$i$ window and 5~s in the Day~$i+1$ session to be considered in the analysis. Correlations were first calculated within recording cycle and then averaged with equal weight across animals for group curves.

Correlation was also measured across a sliding pivot between equal, abutting 30 min windows. Pivots were advanced by 60~s. Within each recording cycle, Day~$i$ and Day~$i+1$ were restricted to pivot positions valid in both sessions, beginning at 30~min and ending at that cycle's last shared pivot. Each window side required at least 300~s of included calcium data. 

\subsubsection*{Cross session position decoding}

For each session pair, a Poisson naive-Bayes decoder was trained on one session and applied to the other \cite{zhang_interpreting_1998}. Place maps in the training session defined the expected activity $\lambda_{jn}$ of neuron $n$ at spatial location $j$. The Poisson likelihood was used for non negative calcium activity; calcium amplitudes were not interpreted as literal spike counts.

For an activity vector $\mathbf q_t$, position was decoded as
\begin{center}
\(\widehat j_t=\arg\max_j\left[
\sum_n q_{tn}\log(\lambda_{jn})
-\Delta t_{\mathrm{eff}}\sum_n\lambda_{jn}
+\log(\pi_j)
\right],
\)
\end{center}
where $\mathbf q_t$ was the vector of calcium activity integrated over the decoding window, $\pi_j$ was the occupancy prior and $\Delta t_{\mathrm{eff}}=0.36$~s was the integration duration. Rates and prior probabilities were floored at $10^{-12}$ before taking logarithms. Error was the Euclidean distance on the $15\times15$ maze grid between decoded and observed center of mass locations, expressed in maze-grid units and averaged over test frames.

For the neuron ID shuffle, registered neuron IDs in Day~$i+90$ were permuted relative to the fixed Day~$i$ place maps. The same neurons and target frames were retained, and the decoder was not refitted. Error was averaged across 1,000 shuffles within each session pair. 

For an individual error map, errors were averaged within observed location, and locations with fewer than five test frames were left blank. For the cohort spatial error map, location specific errors were averaged within session pair, within animal and across animals, giving each animal equal weight. Only locations represented in every contributing animal entered the error map.

\subsubsection*{Reward code stability}

Reward decoding error was evaluated at the tile immediately next to the reward end node and at three nonreward tiles matched for tree depth and local geometry. Analyses using Day~$i$ as the source were restricted to pairs in which both sessions were rewarded; pairs with an unrewarded source session were excluded; All otherwise eligible Day~$i+1$ source sessions met this criterion. The primary analysis used every eligible test frame and the empirical training session occupancy prior. 

For decoding analysis with balanced sampling, locations required at least 20 retained frames in both source and target sessions. Pairs required the reward location and at least one qualifying comparison control location. For frame balancing, each draw sampled source and target sessions independently to their respective smallest frame counts across retained locations. A visit comprised consecutive retained samples from the same location separated by no more than five acquisition frames. For visit balancing, source activity and target decoding error were first averaged within visits. Each session was then independently sampled to its smallest visit count across retained locations. Source tuning maps were re-estimated in each draw from sampled frames or equally weighted visit means. Both procedures used a uniform training prior and 1,000 resamples. Resamples were averaged within session pair and then within animal.

Reward location population vector recurrence used Pearson correlations between population vector over all retained samples at each location. Matched reward recurrence compared the reward target with the reward source. Matched control recurrence averaged the three correlations pairing each comparison target with its corresponding comparison source. The cross location control comprised the mean correlation of the reward target with the three comparison sources and the mean correlation of the three comparison targets with the reward source. When each neuron's maze map was first z-scored across locations, the reward advantage was larger (0.18 and 0.15 for Day~$i$ and Day~$i+1$ sources). 

\subsubsection*{Place maps along linearized routes}

Home to water runs followed the unique direct route from the home to the water port; water to home runs followed the reverse route. Only moving frames (nose speed 2 to 60~cm~s$^{-1}$) were used. Approach included moving frames from leaving home through the first 0.5~s from entry into the reward end node, including entry. Departure included moving frames in the last 0.5~s before exit from that node, excluding exit, followed by the run home beginning at exit. Both windows were bounded by the arrival and exit of that visit. A gap longer than 0.4~s between moving frames ended a window, checked forward from entry for approach and backward from the first moving frame at or after exit for departure. Any frame claimed by opposite-direction runs was assigned by the direction of nose movement along the route (faster than 1~cm~s$^{-1}$, estimated from neighbouring moving frames up to 120~ms away on each side); frames without clear direction entered neither map. On the home to water route no frames were claimed by both directions. Ending the windows instead at the first frame slower than 2~cm~s$^{-1}$ gave reward end advantages ranging from 0.08 to 0.14 in the two directions, using the same neurons and physical bins. A direction was analyzed when both sessions contained at least five runs. The full 83.6~cm route from home to the wall holding the water port was divided into 22 equal 3.8~cm bins, with identical physical boundaries for both directions and at least 0.2~s occupancy required in every bin.

Route maps were calculated from deconvolved activity to localize activity more closely to when the mouse passed each position. The slower decay retained by denoised calcium can spread activity into following positions during a directional run and carry activity from before the run into its first positions. Maps were not smoothed. For heatmap display only, each map was divided by its own maximum. All matched neurons were considered, and a neuron was included when it was active on the route in both sessions and at least half of the route bins met the occupancy criterion.

Each neuron's route map was standardized within each session by subtracting its mean across bins and dividing by its standard deviation across bins. Spatial bins were weighted equally, using shared observed bins where needed, rather than weighting by occupancy or time. This preserved the map's peak and shape while removing its average level and overall scale. At each bin, Day~$i$ and Day~$i+90$ population vectors were correlated across matched neurons using Pearson correlation, comparing relative spatial preferences rather than absolute activity. For the same neurons and observed bins, this correlation is unchanged if any neuron's whole map is shifted or multiplied by a positive factor in either session, as is its place-map correlation. The reward end advantage was the correlation at the bin at the water port minus the mean correlation across the other bins. With denoised calcium instead of deconvolved activity, the advantage remained significant in both directions and in the within-day reliability, equal-frame and shorter-route controls. For the neuron ID shuffle, Day~$i+90$ neuron IDs were permuted within each session pair while all maps were held fixed (100 permutations).

To test whether the water port was better preserved only because it was measured more reliably, each session's runs were split into odd and even runs in time order and maps were recalculated from each half. Before standardization, all four half-session maps were restricted to the same spatial bins meeting the occupancy criterion in every half. Within-day similarity was the geometric mean of the two sessions' correlations between odd-run and even-run population vectors, and across-gap similarity was the mean of the four correlations between Day~$i$ and Day~$i+90$ half-session population vectors, so that both comparisons used half-session maps. The reliability-adjusted correlation was their ratio, calculated for the bin at the water port and for the mean of the other bins; it is a ratio of correlations and is not bounded by 1. For each route position, recording pairs contributed only when both within-day and all four across-gap correlations were defined, so both parts of the ratio used the same pairs. Ratios required positive averaged within-day correlations in both sessions; port and elsewhere summaries used only animals with both ratios defined. For the equal frame control, each session's sample size was its smallest occupancy across the 22 bins. That number of frames was drawn from every bin in each of 100 draws, and population vector correlations (with five shuffles per draw) were recomputed before averaging within session pair. A neuron without activity in a half-session map or in a draw was kept, with a map of zeros.

The two shorter route controls began on junctions at different distances from the entrance (nodes 2 and 3) and converged on the water port. Runs toward and away from the port on these routes used the same frame rule at the port, with velocity measured along each route. The node 2 and node 3 routes contained 19 and 15 bins, respectively, with identical boundaries in both directions. Bins were 3.8~cm wide and anchored at the port wall, with a 1.9~cm remainder at the junction end; every bin required at least 0.2~s occupancy. Every included session contained at least five runs in each direction. These shorter-route controls retained the published requirement that reward was available in the earlier session. For the reward delivery control, home to water runs were classified as rewarded when water was delivered within 3~s of arrival in the reward end node and as timeout otherwise, and maps were calculated separately from each run type. A later delivery did not change this classification and lay outside the 0.5~s approach window. Frames were not clipped to prompt delivery, so the rewarded maps included a small number of frames at or after water delivery. Neurons were included when they were active on the route in both sessions on both run types, with at least half of the bins meeting the occupancy criterion, and session pairs when they contained at least ten such neurons and at least five runs of each type in both sessions.

The second day of each visit (Day~$i+1$ and Day~$i+91$) was analyzed in the same way, using the same session pairs when both second-day sessions contained at least five runs in the analyzed direction. Calcium traces from the two days of a visit had been extracted together, so neuron indices were shared between them and the neuron matching between Day~$i$ and Day~$i+90$ also applied to Day~$i+1$ and Day~$i+91$; the stored spatial footprints of the two days were identical for all 45 visits examined. Neurons were selected on the second-day maps, using only spatial bins meeting the occupancy requirement in both compared recordings for normalization and correlation; unobserved bins were excluded. For the comparison of approach and departure, session pairs analyzed in both directions were used (27 pairs from 13 animals), with matched neurons active on both routes in both sessions. Water to home maps were ordered from the home to the water port, Day~$i$ approach maps were correlated with Day~$i+90$ departure maps and Day~$i$ departure maps with Day~$i+90$ approach maps, and the two correlations were averaged.

For route analyses, session pair values were averaged within animal, and animals were weighted equally; for the reliability-adjusted correlation, the within-day and across-gap correlations were first averaged over an animal's session pairs (and, for the rest of the route, over its bins), and the ratio was then formed for each animal. Confidence intervals were Student's $t$ intervals across animals. Differences were tested with exact two sided sign-flip tests, which compare the observed mean across the $G$ animals with the means obtained under all $2^G$ assignments of signs to the animal values. BH correction was applied across the two directions in Fig.~\ref{fig:figure4}e, across the eight comparisons between the water port and the rest of the route in Extended Data Fig.~\ref{fig:ed08}e,f,h, across the two directions for the second-day analysis, and across the two run types in Extended Data Fig.~\ref{fig:ed08}i. The difference between directions, the rewarded minus timeout difference and the comparison of approach and departure were single tests. 

\subsubsection*{Level maps and controls}

Two related partitions were used and are named separately throughout. The primary level maps comprised seven levels (0--6) corresponding to tree depth. The geometric control comprised eight levels by a different assignment of the junction tiles. 

The level decoder used the same Poisson model with level maps and the empirical occupancy prior across the seven primary tree levels. Calcium activity was integrated over 0.36~s window, and each training level required 40~s of occupancy. Accuracy was the fraction of test frames assigned to the correct level. For the neuron ID shuffle, registered neuron IDs in Day~$i+90$ activity were permuted relative to the fixed Day~$i$ level maps without refitting the decoder. All eligible target frames were decoded, and accuracy was averaged across 200 shuffles within each session pair. Level decoding was repeated after excluding neurons whose place map correlation exceeded thresholds from 0.1 to 0.7.

A neuron's preferred level was the unique maximum of its level map. Maps that were absent, constant, tied at the maximum or insufficiently occupied had no preference. Among neurons whose preferred level differed between sessions, transition probabilities were normalized within each source level. The null distribution was obtained by permuting destination preferences within session pair 10,000 times and identifying neurons with changed preferences anew after each permutation. 

To disrupt correspondence between exact locations across sessions while preserving level maps, later session place maps were randomly permuted 1000 times among all eligible connected components assigned to the same level. The same spatial rearrangement was applied to every neuron. Smoothing was then applied separately to each neuron's transformed map to restore the continuity of the original place maps.

A matched resolution control removed the narrow connecting corridors, divided the remaining maze into sixteen square blocks and merged adjacent block pairs along either the horizontal or vertical axis to form eight bins. Geometric, horizontal and vertical maps were calculated from calcium activity using the registration specific to each comparison. Each bin required at least 5~s of data in each session. 

\subsubsection*{Neural--behavior associations}

\paragraph*{Primary associations.}
Reward decoding error from the rewarded Day i source session was regressed against first long bout time in the later session. Neurons were registered across Day~$i$, the paired Day~$i+1$ source and the later target session. Ensemble size was set by the smallest qualifying shared population across the included session pairs, allowing every comparison to use an equally sized population. Neurons were sampled in 50 random draws, while one fixed 80~s sample of reward location activity from the target session was held constant. Each decoder retained the empirical source session occupancy prior. For the primary structural code association, the mean level correlation from each successive long gap (3 month) Day~$i$ session pair was regressed against $E_{32}$ in the later session. Each session pair contributed one value. 

\paragraph*{Cross measure comparison from matched population.}
For cross measure comparisons, analysis was restricted to successive long gap session pairs in which all six neural readouts and the relevant later session behavioral measures were defined. Within this common set, reward decoding error, whole-maze decoding error, place map correlation, level map correlation, population vector correlation and ensemble rate correlation were recalculated from coordinated, equally sized samples of registered neurons. Population size was set by the smallest qualifying registered population across these pairs. Neurons were sampled in 50 random draws, and the same identities were used for all six measurements within a draw. Reward decoding used 80~s of test activity from the reward location in every pair. Whole-maze decoding used 80~s sampled across the maze without balancing occupancy. The number of spatial locations was set by the session pair with the fewest source--target locations in common; each pair sampled this number from its own common locations. For the matched population level maps, location means were pooled over levels 0--5; level~6 was excluded because it was not represented in 3 included session pairs. For ensemble rate, 80,000 frames were equally sampled from each session. 

\paragraph*{Commonality analysis.}
Commonality analysis separated the behavioral variance associated with reward decoding error and level map correlation into components unique to each neural measure and a component shared by both. Each unique component was the additional variance explained after accounting for the other neural measure. The shared component represented variance explained by both measures. We tested whether reward decoding error and level map correlation had different relationships with the two behaviors. Decoding error and first long bout time were sign-reversed so that higher values consistently indicated better decoding, earlier bout onset or better performance. All variables were standardized, and each session pair received equal weight in a stacked linear model. The model compared the two neural measures for each behavior and then tested whether this difference changed between behaviors.

\subsection*{Statistics and reproducibility}

Analyses were performed in Python~3.10.9 using NumPy~1.26.4, pandas~2.3.3, SciPy~1.10.0 and statsmodels~0.13.5. Statistical plots were generated with Matplotlib~3.7.0.

Unless stated otherwise, linear models accounted for repeated observations from the same animal by clustering standard errors by animal. Standard errors were adjusted for the finite number of animals, and tests used a Student's $t$ distribution with $G-1$ degrees of freedom, where $G$ was the number of animals. Paired comparisons tested whether the mean within-pair difference differed from zero. When an animal contributed multiple pairs, these differences were clustered by animal. Bootstrap confidence intervals were obtained by sampling animals with replacement while keeping all observations from each sampled animal together. Percentile intervals were used unless stated otherwise. Tests were two-sided unless a directional test was specified. When related tests formed a correction family, their $p$ values were corrected using the method specified for that analysis. Benjamini--Hochberg correction was used to control the false discovery rate. Figure asterisks denote $p<0.05$ (*), $p<0.01$ (**), and $p<0.001$ (***). They represent corrected $p$ values for family-controlled tests and raw $p$ values for standalone tests. 

%\newpage
\section*{Acknowledgements}

We thank Walter Gonzalez for help setting up the custom miniscope technique; Matthew Rosenberg for maze construction; Nicole Schweers for animal support; Daniel Wagenaar of the Caltech Neurotechnology Lab for engineering assistance; the Caltech HPC for computing resources; Thanos Siapas, xxx for critical feedback; and all Meister lab and Lois lab members for their input. This work was supported by grants from the National Institutes of Health (NIH) (5R01NS111477 to M.M. and 5R01NS121918 to C.L.) and the Simons Collaboration on Global Brain (543015 to M.M.).  

\section*{Author contributions}

J.W., M.M. and C.L. conceived the study,  J.W. conducted the experiments and analyzed the data.  J.W. wrote the paper with input from M.M. and C.L. M.M. and C.L. supervised all aspects of the work. 

\section*{Competing interests}

The authors declare no competing interests. 

\section*{Data and code availability}

Data and code will be made available in online repositories upon publication. 

\newpage
\bibliography{references}

\clearpage
\nolinenumbers
\begin{center}
\saveimageanchor{fig:figure1}
\includegraphics[
    width=\textwidth,
    height=0.95\textheight,
    keepaspectratio
]{figures/figure1_version_record.pdf}

\end{center}

\clearpage

\captionof{figure}{
\textbf{Lifelong memory of a complex navigation task.} \\
\textbf{a,} Schematic of the binary-tree maze, comprising six branch decisions and 64 end nodes in a $15\times15$ grid. The home cage entrance and fixed water port are indicated by the home and water drop icons, respectively. Colored terminal regions define the end-node groups used in the end node timelines in \textbf{c} and \textbf{f}. The red path illustrates an example long bout, defined as a sustained approach to the rewarded end node containing more than six consecutive correct decisions. Circles numbered 0 to 5 mark the six decision nodes scored in \textbf{h}.\\
\textbf{b,} Longitudinal design of recurring two-day recording cycles. Each cycle comprised two sessions of behavioral and neural recording (Day $i$ and Day $i+1$), followed by approximately 90 days in the home cage before the next cycle. The Day $i$ session was naive in the first cycle and a return session thereafter.\\
\textbf{c,} Example end node visit timelines on Day 1, Day 2, Day 91 and Day 181. Each point marks an end node visit, colored by the terminal region in \textbf{a}; the horizontal reference line and water drop indicate the rewarded end node. In this example, compared with naive Day 1, visits on Day 2, 91 or 181 were less clustered from the beginning and concentrated near the reward end node earlier in the session, indicating efficient navigation and rapid recovery of runs towards reward.\\
\textbf{d,} First long bout time across repeated 3 month cycles. Naive mice required more time in the maze before the first long bout than mice returning after months away. First long bout time fell from $1{,}980 \pm 266$ s on the naive day to $650 \pm 127$, $446 \pm 74$, $217 \pm 43$ and $259 \pm 51$ s after 3, 6, 9 and 12 months, respectively (mean $\pm$ s.e.m.; $n=23$, 17, 17, 12 and 10 sessions, respectively). Each return group differed from naive (regression on log time, with uncertainty adjusted for repeated sessions from the same animal; Benjamini--Hochberg (BH) corrected $p$ values were $9.7\times10^{-4}$, $2.1\times10^{-6}$, $1.5\times10^{-6}$ and $2.0\times10^{-5}$ at 3, 6, 9 and 12 months). Across return sessions, first long bout time declined with each additional recording cycle ($n=56$ sessions from 17 animals; fold change per return cycle, 0.71; two-sided $p=0.036$). Points denote sessions; bars and error bars show mean $\pm$ s.e.m.\\
\textbf{e,} Distribution of exploration efficiency for naive and return sessions. Return sessions showed less redundant end node sampling than naive sessions. Exploration efficiency is defined as $32/N_{32}$, with $N_{32}$ denoting the number of end node visits required to sample the first 32 distinct end nodes. Efficiency increased from $0.12\pm0.01$ in naive sessions to $0.37\pm0.02$ in return sessions (mean $\pm$ s.e.m.; 23 naive sessions from 23 animals and 56 return sessions from 17 animals). Histograms show the fraction of sessions within each condition. \\
\textbf{f,} Example end node visit timelines from latent control. The water port was replaced by a maze wall on latent day1 and restored on latent day2. Point colors follow the terminal end node groups in \textbf{a}. At the beginning of latent day1, latently trained mice showed an inefficient exploration pattern similar to that of animals on naive days with reward. They were efficient at the beginning of latent day2.\\
\textbf{g,} First long bout time in naive, latent day2 and return sessions. Animals in latent day2 sessions had recent layout familiarity but no prior reward experience, and reached the first long bout faster than naive animals ($706\pm139$ s versus $1{,}980\pm266$ s). Return animals were faster still ($425\pm51$ s; mean $\pm$ s.e.m.; $n=23,9$ and 56 sessions from 23, 9 and 17 animals). BH corrected $p$ values were 0.0068 for latent day2 versus naive, $5.9\times10^{-8}$ for return versus naive and 0.016 for return versus latent day2.\\
\textbf{h,} First visit route score, defined as the fraction of the six reward path decision nodes at which the animal’s first choice was directed towards reward. Scores are shown for naive, return, second day and latent day2 sessions ($n=23,56,74$ and 9 sessions from 23, 17, 24 and 9 animals). Second day scores exceeded naive, return and latent day2 scores (BH corrected $p=4.4\times10^{-4}$, $1.6\times10^{-4}$ and $0.044$). The other comparisons were not significantly different after correction (return versus naive, corrected $p=0.16$; latent day2 versus naive, corrected $p=0.57$; return versus latent day2, corrected $p=0.74$). Points denote sessions; bars show condition means with 95\% confidence intervals from regressions with standard errors clustered by animal. The dashed horizontal line marks chance performance (0.5).
}
\labeltoimage{fig:figure1}

\clearpage



\nolinenumbers
\begin{center}
\saveimageanchor{fig:figure2}
\includegraphics[
    width=\textwidth,
    height=0.95\textheight,
    keepaspectratio
]{figures/figure2_version_record.pdf}
\end{center}

\clearpage

\captionof{figure}{
\textbf{CA1 retains a persistent trace after long periods without maze experience.} \\
\textbf{a,} Example longitudinal CA1 local correlation images from Day 1 and Day 91, showing a stable field of view. Scale bar, $100\,\mu\mathrm{m}$. \\
\textbf{b,} Four example place maps from matched neurons on Day $i$ and Day $i+90$. The examples span retained and remapped spatial fields, with place map correlations of 0.63, 0.42, 0.18 and $-0.20$. Each red arrow represents a calcium event, and the underlying gray lines show animal trajectories.\\
\textbf{c,h,m,} Schematics defining place map, population vector and ensemble rate correlations, respectively. Place map correlation compares a neuron's rate vector normalized by position between sessions; population vector correlation compares neuron activity vectors at corresponding positions; ensemble rate correlation compares vectors of each matched neuron's mean activity rate between sessions. \\
\textbf{d,i,n,} Distributions of place map, population vector and ensemble rate correlations, respectively, for Day $i$ to Day $i+90$ comparisons. The place map distribution was centered near zero but carried a high correlation tail; population vector correlations showed a broader positive shift; ensemble rate correlations were strongly positive across long gaps. Neuron ID shuffle controls are shown alongside the data. The analyses comprised 36 session pairs from 14 animals, yielding $16,901$ matched neuron observations with sufficient calcium events, $3,810$ spatial bin observations and one ensemble rate correlation per session pair. \\
\textbf{e,j,o,} Average place map, population vector and ensemble rate correlations across a 3 month gap compared with 10,000 neuron ID shuffles. Observed correlations were 0.12, 0.22 and 0.53 across session pairs, respectively, compared with shuffle means of 0.0028, $-3.0\times10^{-6}$ and $-5.4\times10^{-5}$ (one-sided permutation $p=1.0\times10^{-4}$ for each metric). \\
\textbf{f,k,p,} Place map, population vector and ensemble rate correlations across intervals from Day 1 to Day 2 through 12 months. All three correlations declined from adjacent days to month scale gaps but remained above shuffle through 9 months. Place map correlations at 1, 90, 180, 270 and 360 days averaged across session pairs were 0.53, 0.12, 0.08, 0.13 and 0.01. The corresponding population vector correlations were 0.60, 0.22, 0.17, 0.18 and 0.05, and ensemble rate correlations were 0.73, 0.53, 0.42, 0.35 and 0.20. Values are means across session pairs with 95\% confidence intervals. Gray dashed lines show neuron ID shuffles. The 1 day point summarizes 47 Day 1 to Day 2 session pairs within the same cycle from 14 animals. The 90, 180, 270 and 360 day points summarize 36, 22, 11 and 4 session pairs comparing Day 1 sessions from successive recording cycles, from 14, 11, 7 and 4 animals, respectively. \\
\textbf{g,l,q,} Recurrence of place map, population vector and ensemble rate correlations across successive long gaps. Higher correlation in one gap was associated with higher correlation in the next. In \textbf{g,l}, previous gap correlations were divided into equal width bins across the x axis. The solid curves show the mean next gap correlation at each bin center. Faint points show the underlying 9,662 matched neuron and 2,266 spatial bin observations, respectively. In \textbf{q}, points show 22 three session chains and Pearson $r=0.71$. The line and shading show the fitted regression and its animal-clustered 95\% confidence interval. Each analysis included 11 animals with at least three recorded cycles. Regressions of next gap on previous gap correlation, with standard errors clustered by animal, gave slopes of 0.36, 0.43 and 0.35 for place maps, population vectors and ensemble rates. The corresponding BH corrected $p$ values were $7.8\times10^{-4}$, $7.8\times10^{-4}$ and 0.028.
}
\labeltoimage{fig:figure2}

\clearpage

\nolinenumbers
\begin{center}
\saveimageanchor{fig:figure3}
\includegraphics[
    width=\textwidth,
    height=0.95\textheight,
    keepaspectratio
]{figures/figure3_version_record.pdf}
\end{center}

\clearpage

\captionof{figure}{
\textbf{Return sessions rebuild detailed maps from an experience-dependent initial state.} \\
\textbf{a,} Example place maps from one neuron across two successive recording cycles. Spatial structure recurred between cycles, but the maps also changed between the first and second session within each cycle. \\
\textbf{b,} Analysis schematic showing moving 30 min windows from Day $i$ compared with the full Day $i+1$ representation used as the template.\\
\textbf{c,} Representations from the first 30 min of Day $i$ were compared with the full Day $i+1$ template across recording cycles.\\
\textbf{d,} Cross cycle comparison of Day $i$ representations from successive recording cycles.\\
\textbf{e,h,k,} Place map, population vector and ensemble rate correlations, respectively, across moving 30 min windows on Day $i$. Correlations generally increased with time in the maze and began higher after prior experience. Curves show means across animals $\pm$ s.e.m. Across \textbf{e--m}, increasing color opacity indicates greater cumulative experience.\\
\textbf{f,i,l,} Early alignment increased with experience ($n=59$ observations from 22 animals; Pearson $r=0.62$, 0.51 and 0.39; BH corrected $p=5.4\times10^{-7}$, $1.2\times10^{-4}$ and 0.0016). Points represent individual animals in individual recording cycles; lines are ordinary least squares fits and shading shows 95\% confidence intervals.\\
\textbf{g,j,m,} Long gap correlation across successive recording cycles ($n=37$ session pairs from 14 animals, including one short interval pair). Place map and population vector correlations increased across returns (Pearson $r=0.49$ and 0.43; BH corrected $p=0.037$ for each), whereas ensemble rate correlation showed no supported trend ($r=0.20$, corrected $p=0.14$).
}
\labeltoimage{fig:figure3}

\clearpage

\nolinenumbers
\begin{center}
\saveimageanchor{fig:figure4}
\includegraphics[
    width=\textwidth,
    height=0.95\textheight,
    keepaspectratio
]{figures/figure4_version_record.pdf}
\end{center}

%\clearpage

\captionof{figure}{
\textbf{The reward location anchors selective place stability.} \\
\textbf{a,} Schematic of Bayesian spatial decoding across sessions. Place maps from Day $i$ were combined with calcium activity from matched neurons on Day $i+90$ to decode position across a long gap. The gray region indicates the 360 ms time window used for decoding. \\
\textbf{b,} Cross session spatial decoding error compared with neuron ID shuffle controls. Decoded position was slightly more accurate than after neuron ID shuffling (data, $5.7\pm 0.2$ bins; neuron ID shuffle, $6.6\pm 0.1$ bins; mean $\pm$ s.e.m.; $n=36$ session pairs from 14 animals; two sided regression with standard errors clustered by animal, $p=0.014$). Lines show paired observations. \\
\textbf{c,} Example maze map of spatial decoding error for one 0 to 3 mo session pair. Decoding error was low near the reward location (2.0 bins) and higher at matched control locations (7.0, 7.0 and 8.2 bins). The orange outlined decoder bin is the reward bin, and the gray outlined bins are matched controls. \\
\textbf{d,} Reward decoding error compared with matched control locations across session pairs whose source Day $i$ was rewarded. Error was lower at the decoder bin adjacent to the water port ($3.7\pm 0.5$ bins) than at three matched controls ($7.0\pm 0.3$ bins). The analysis included 30 session pairs from 14 animals (two sided regression with standard errors clustered by animal, $p=1.2\times10^{-4}$).\\
\textbf{e,} Linearized place maps along home to water (left) and water to home (right) runs for matched neurons on Day $i$ and Day $i+90$, with frames at the water port taken from the first 0.5~s after arrival and the last 0.5~s before departure (Methods). Each run reads from left to right, and bins are numbered from the home (0) to the water port (21) in both directions. Each map is scaled to its maximum for display, with the colour scale saturating at 0.6, and rows in both sessions are ordered by the Day $i$ peak location, with neurons peaking at the water port at the bottom. Line plots show the population vector correlation between Day $i$ and Day $i+90$ at each route position after each neuron's map was standardized (mean across animals; shading, 95\% CI), with the mean of 100 neuron ID shuffles (dotted). In both directions the correlation peaked at the water port, where it exceeded the mean of the other positions by 0.14 on home to water runs (95\% CI, 0.10 to 0.19; BH corrected $p=4.9\times10^{-4}$) and by 0.11 on water to home runs (0.05 to 0.17; corrected $p=0.0023$; two sided exact sign-flip tests across animals). The home to water analysis included 8,222 matched neuron observations from 29 session pairs and 13 animals; the water to home analysis included 9,017 observations from 28 session pairs and 14 animals. \\
\textbf{f,} Reward decoding error plotted against first long bout time in the later session. Session pairs with lower reward decoding error reached the first long bout sooner ($n=30$ session pairs whose source Day $i$ was rewarded, from 14 animals; Pearson $r=0.46$; two sided regression with standard errors clustered by animal, $p=0.0041$). 
}
\labeltoimage{fig:figure4}

\clearpage

\nolinenumbers
\begin{center}
\saveimageanchor{fig:figure5}
\includegraphics[
    width=\textwidth,
    height=0.95\textheight,
    keepaspectratio
]{figures/figure5_version_record.pdf}
\end{center}

%\clearpage

\captionof{figure}{
\textbf{A compact structural code predicts memory guided exploration.} \\
\textbf{a,} Schematic of the tree level representation. Maze locations were grouped by depth in the binary-tree hierarchy into seven levels, independent of exact two dimensional position.\\
\textbf{b,} Heatmap of level tuning across 16,901 matched neuron observations from 36 long gap session pairs and 14 animals. \\
\textbf{c,} Distribution of level map correlations across matched neuron observations from long gap session pairs. Correlations were shifted positive and included a peak near one (mean, 0.27; $n=16{,}901$ matched neuron observations).\\
\textbf{d,} Average level map correlation compared with 10,000 neuron ID shuffle controls. Level map stability exceeded shuffle (data, 0.25; shuffle, 0.035; one sided permutation $p=1.0\times10^{-4}$).\\
\textbf{e,} Cross session level decoding accuracy compared with neuron ID shuffle controls. Tree depth was decoded above shuffle (data, $0.38\pm0.01$; neuron ID shuffle, $0.24\pm0.01$; mean $\pm$ s.e.m.; $n=36$ session pairs from 14 animals; one sided regression with standard errors clustered by animal, $p=9.2\times10^{-7}$).\\
\textbf{f,} Level map correlation plotted against exploration efficiency. Session pairs with stronger preserved level tuning sampled end nodes more efficiently ($n=36$ session pairs from 14 animals; Pearson $r=0.75$; two sided regression with standard errors clustered by animal, $p=1.9\times10^{-4}$). 

}
\labeltoimage{fig:figure5}

\clearpage

\nolinenumbers
\begin{center}
\saveimageanchor{fig:ed01}
\includegraphics[
    width=\textwidth,
    height=0.95\textheight,
    keepaspectratio
]{figures/ed01_version_record.pdf}
\end{center}

\clearpage
\captionof{extendeddatafigure}{
\textbf{Additional reward and exploration measures for lifelong spatial memory.}\\
\textbf{a,} Trajectories matching the examples in Fig.~\ref{fig:figure1}c from the start of each session until the first reward visit on Day 1, Day 2, Day 91 and Day 181. The shared color scale spans 0 to 16 min in maze; the home entrance and water port are indicated by the home and water drop icons.\\
\textbf{b,} Cumulative rewards as a function of time in maze on the naive days, second days and return days after 3, 6, 9 and 12 months. Curves show session means $\pm$ s.e.m. through the last time point with at least 80\% of sessions in each group (23, 74, 17, 17, 12 and 10 sessions, respectively).\\
\textbf{c,} Cumulative long bouts over time in maze for the same groups as in \textbf{b}.\\
\textbf{d,} Cumulative rewards as a function of maze tile transitions. Curves show means $\pm$ s.e.m. over the transition range with at least 80\% of sessions as in \textbf{b}.\\
\textbf{e,} Cumulative long bouts as a function of maze tile transitions for the same groups as in \textbf{b}.\\
\textbf{f,} Time in maze and maze tile transitions until the first reward, tenth reward and first long bout. Points are sessions and bars show mean $\pm$ s.e.m. for naive day, return day and second day sessions (23, 56 and 74 sessions, respectively). Outcomes were analyzed after $\log(1+x)$ transformation using linear models with standard errors clustered by animal; all displayed comparisons were significant after BH correction (maximum corrected $p=1.7\times10^{-4}$).\\
\textbf{g,} Maze tile transitions plotted against time in maze before the first long bout. The slope between the log transformed variables was near unity (slope, 0.964; 95\% CI, 0.94 to 0.99; two sided regression with standard errors clustered by animal, $p=2.7\times10^{-35}$; 153 sessions from 29 animals). Of these animals, 23 contributed naive day sessions. The other six had initial sessions that were latently learned, so they contributed only later return day and second day sessions to this analysis. Thus, elapsed time and traversed distance accumulated approximately proportionally across sessions.\\
\textbf{h,} Center of mass tracking speed during exploration and reward directed modes within the first 30 min of each session. Symbols and error bars show mean $\pm$ s.e.m. on Day 1 and Day 2 at 0, 3, 6, 9 and 12 months. For exploration, the respective sample sizes were 23, 17, 17, 12 and 10 Day 1 sessions and 18, 17, 17, 12 and 10 Day 2 sessions. For the reward directed mode, they were 20, 17, 17, 12 and 10 sessions on Day 1 and 18, 17, 17, 12 and 10 sessions on Day 2.
\textbf{i,} Illustrative repeated coverage rounds from example naive and return sessions. End node visits are shown at left with distinct end nodes colored within each round. Middle plots show the accumulation of 32 distinct end nodes within successive coverage rounds; dashed lines mark coverage round boundaries. Right plots show exploration efficiency across coverage rounds, with filled symbols indicating the exploration efficiency used in the main text.\\
\textbf{j,} Exploration efficiency as a function of time in maze for naive day, return day and second day sessions as in \textbf{b}. Curves show mean $\pm$ s.e.m. in fixed 10 min bins; only bins with at least five sessions are shown.\\
\textbf{k,} Exploration efficiency as a function of maze tile transitions, summarized as in \textbf{j} in fixed bins of 1,000 transitions.

}
\labeltoimage{fig:ed01}


\clearpage

\nolinenumbers
\begin{center}
\saveimageanchor{fig:ed02}
\includegraphics[
    width=\textwidth,
    height=0.95\textheight,
    keepaspectratio
]{figures/ed02_version_record.pdf}
\end{center}

\clearpage
\captionof{extendeddatafigure}{
\textbf{Junction turning policies and biased walk models of exploration efficiency.}\\
\textbf{a,} Four local turning policies at junctions. $\mathrm{P}_{\mathrm{SF}}$ is the probability of continuing from the stem into a branch rather than reversing. Conditional on continuing, $\mathrm{P}_{\mathrm{SA}}$ is the probability of choosing the alternating branch. $\mathrm{P}_{\mathrm{BF}}$ is the probability of leaving a previously sampled subtree rather than reentering it. Conditional on leaving, $\mathrm{P}_{\mathrm{BS}}$ is the probability of continuing towards the stem rather than entering the sibling branch.\\
\textbf{b,} Policies estimated for each session on naive and return days, restricted to sessions in which all four policies were estimable (23 sessions from 23 animals and 58 sessions from 18 animals, respectively). Hollow symbols are sessions, solid symbols are group means and dashed lines are equal choice references. Return minus naive differences were $+0.01$ for $\mathrm{P}_{\mathrm{SF}}$ (BH corrected $p=0.38$) and $+0.09$ for $\mathrm{P}_{\mathrm{BF}}$ (corrected $p=4.6\times10^{-9}$). The corresponding differences were $+0.12$ for $\mathrm{P}_{\mathrm{SA}}$ (corrected $p=9.1\times10^{-9}$) and $+0.09$ for $\mathrm{P}_{\mathrm{BS}}$ (corrected $p=1.9\times10^{-10}$). Policies were estimated from exploration choices during the first 30 min in the maze. \\
\textbf{c,} Return day and second day policies, restricted to sessions in which all four policies were estimable (58 sessions from 18 animals and 74 sessions from 24 animals). Differences between second day and return day were $-0.003$ for $\mathrm{P}_{\mathrm{SF}}$ (BH corrected $p=0.53$) and $-0.04$ for $\mathrm{P}_{\mathrm{BF}}$ (corrected $p=4.0\times10^{-7}$). The corresponding differences were $-0.05$ for $\mathrm{P}_{\mathrm{SA}}$ (corrected $p=1.5\times10^{-5}$) and $+0.02$ for $\mathrm{P}_{\mathrm{BS}}$ (corrected $p=3.4\times10^{-4}$).\\
\textbf{d,} Contribution of each return day policy to simulated exploration efficiency, calculated as efficiency before replacement minus efficiency after replacing that policy with the same animal's naive value while holding the other policies fixed (58 return sessions from 18 animals). The contributions of $P_{\mathrm{SF}}$ and $P_{\mathrm{BF}}$ were $-0.4$ percentage points (95\% CI, $-1.0$ to $+0.2$) and $+6.1$ points ($+4.4$ to $+8.2$), respectively. The corresponding contributions of $P_{\mathrm{SA}}$ and $P_{\mathrm{BS}}$ were $-2.6$ points ($-3.2$ to $-2.0$) and $+7.2$ points ($+6.0$ to $+8.4$). Alternating left and right choices imposes a local pattern without necessarily moving the animal out of a sampled subtree, whereas efficient coverage in this maze depends on leaving sampled regions and returning toward branch points that open access to larger unsampled subtrees. Positive values indicate greater efficiency with the observed return day policy. \\
\textbf{e,} Absolute change of naive day or second day policy mean from the return day mean. Second day values were closer to return day values for all four policies. \\
\textbf{f,} Observed exploration efficiency versus efficiency generated by the four policy biased walk (23 naive day, 58 return day and 74 second day sessions). Circles denote sessions and dashed diagonals mark identity. Pearson $r$ and mean bias were descriptive measures of session-level agreement (naive day, $r=0.11$, bias $=+0.16$; return day, $r=0.30$, bias $=+0.05$; second day, $r=0.06$, bias $=+0.02$).\\
\textbf{g,} Stemward exit probability $\mathrm{P}_{\mathrm{BS}}$ across shallow, middle and deep non-root decision junctions. These included all junctions at the same tree depths corresponding to reward route decision nodes 1–2, 3–4 and 5, respectively, in Fig.~\ref{fig:figure1}a. The analysis used 22 naive day sessions from 22 animals, 58 return day sessions from 18 animals and 74 second day sessions from 24 animals with sufficient samples at all three depths. Upper curves show session means with interquartile bands; lower estimates show group contrasts with 95\% CIs. BH corrected $p$ values for return minus naive differences were $2.5\times10^{-9}$, $2.4\times10^{-13}$ and 0.025 from shallow to deep. The corresponding corrected $p$ values for differences between second day and return day were $3.5\times10^{-11}$, $6.9\times10^{-10}$ and $1.6\times10^{-4}$.\\
\textbf{h,} Reconstruction of the observed return minus naive exploration efficiency gap using the same cohort as in \textbf{g} excluding one naive session without enough samples. The four policy biased walk reconstructed 40.5\% of the observed gap (95\% CI, 28.6 to 54.1\%), whereas allowing $\mathrm{P}_{\mathrm{BS}}$ to vary with depth reconstructed 86.7\% (68.5 to 107.2\%). The common four-policy and depth-varying cohort contained 54 return sessions paired with 17 naive day sessions from 17 animals. Circles and horizontal intervals show reconstructed gaps and 95\% confidence intervals; the vertical black line marks the observed return minus naive gap.\\
\textbf{i,} Observed versus simulated efficiency for the model in which $\mathrm{P}_{\mathrm{BS}}$ varied with depth, using the same all-depth session support and display conventions as in \textbf{g}. Pearson $r$ and mean bias were $r=0.30$ and $+0.032$ for naive day, $r=0.56$ and $+0.030$ for return day, and $r=-0.06$ and $+0.054$ for second day. 
}
\labeltoimage{fig:ed02}


\clearpage

\nolinenumbers
\begin{center}
\saveimageanchor{fig:ed03}
\includegraphics[
    width=\textwidth,
    height=0.95\textheight,
    keepaspectratio
]{figures/ed03_version_record.pdf}
\end{center}

%\clearpage
\captionof{extendeddatafigure}{
\textbf{Reward path fluency and first visit route choices.}\\
\textbf{a,} Mean dwell time at internal junctions during exploration before the first reward encounter in latent day2 and return sessions. Return animals showed shorter junction dwell before reward, consistent with the hypothesis that animals were motivated for an anticipated reward before the reward was found ($\beta_{\mathrm{return-latent}}=-0.31$ s; $\beta$ stands for estimated difference; 55 return day and 8 latent day2 sessions from 20 animals with enough pre-reward samples; BH corrected $p=0.0097$).\\
\textbf{b,} Mean exploratory junction dwell after the first reward. The return minus latent difference was not significant after correction (56 return day and 9 latent day2 sessions from 21 animals; corrected $p=0.066$).\\
\textbf{c,} Difference between dwell at junctions on the direct reward path and other junctions before the first reward. There was no supported difference specific to the reward path before reward ($\beta_{\mathrm{return-latent}}=+0.18$ s; 53 return day and 8 latent day2 sessions from 20 animals; corrected $p=0.15$).\\
\textbf{d,} Junction dwell along the direct home to water path during reward directed approaches after the first reward. Return animals traversed these junctions faster than latent day2 controls ($\beta_{\mathrm{return-latent}}=-0.19$ s; 56 return day and 9 latent day2 sessions from 21 animals; corrected $p=2.6\times10^{-4}$). In \textbf{a--d}, points are sessions, bars are group means and black error bars are 95\% confidence intervals from regressions with standard errors clustered by animal. \\
\textbf{e,} Six decision nodes used for the first visit reward path analysis. The direct home to water path follows the sequence left, right, right, left, right and left; circles numbered 0 to 5 mark the scored junctions.\\
\textbf{f,} Probability that the first choice before reward at each scored junction was directed towards reward for naive day, second day, return day and latent day2 sessions. This negative result still left a possible weaker interpretation: return animals might retain a partial bias toward the reward path even if they no longer preserved the full route sequence. We therefore asked whether the small residual reward-directed bias on return days could be explained by the general branch-choice policies described above. The most relevant policy was alternation, which increased on return days. The home-to-water sequence ($\mathrm{L\,R\,R\,L\,R\,L}$) contains four alternations across five successive turn transitions, so an animal following a non-specific alternation tendency would often choose the reward-directed branch at those nodes without storing the full left--right sequence. Consistent with this explanation, unrewarded latent controls matched return animals at the reward-path choices compatible with alternation. Bars show means across animals $\pm$ s.e.m.; the dashed line marks chance probability 0.5. Sample sizes were 23 sessions from 23 animals, 74 sessions from 24 animals, 56 sessions from 17 animals and 9 sessions from 9 animals, respectively. Second day animals showed reward directed choices across most of the route, whereas the return day profile did not show uniform reinstatement of all six first visit choices. 
}


\labeltoimage{fig:ed03}


\clearpage

\nolinenumbers
\begin{center}
\saveimageanchor{fig:ed04}
\includegraphics[
    width=\textwidth,
    height=0.95\textheight,
    keepaspectratio
]{figures/ed04_version_record.pdf}
\end{center}

\clearpage
\captionof{extendeddatafigure}{
\textbf{Cross session field of view stability and neuron matching quality.}\\
\textbf{a,} Example longitudinal CA1 local correlation images on Day 1, Day 91, Day 181 and Day 271, showing a stable field of view across repeated long gaps. Scale bar, $100\,\mu\mathrm{m}$.\\
\textbf{b,} The same fields of view with cells registered across all four Day 1 sessions shown in green and other detected cells shown in white.\\
\textbf{c,} Cumulative two dimensional field of view displacement across successive recording cycles. Each trajectory shows one animal, with the first field of view set to the origin (14 animals).\\
\textbf{d,} Displacement from the first field of view across recording cycles. Thin lines show individual animals and black points show mean $\pm$ s.e.m. Mean displacement at cycle counts 1 to 4 was 60.9, 84.4, 151.6 and $257.0\,\mu\mathrm{m}$, respectively ($n=13,12,7$ and 4 animals). \\
\textbf{e,} Displacement from the preceding field of view across recording cycles, expressed as a percentage of the preceding field of view diagonal. Mean displacement at cycle counts 1 to 4 was 4.99\%, 3.83\%, 3.60\% and 2.91\%, respectively ($n=13,12,7$ and 4 animals; pooled median, 3.45\%). \\
\textbf{f,} Centroid distance distributions for candidate neuron matches within the $14\,\mu\mathrm{m}$ search radius, evaluated in both directions within each session pair. The nearest candidate by centroid was classified as centroid closest and all remaining candidates as other candidates. The analysis comprised 57,754 centroid closest and 61,690 other candidate observations from 36 session pairs and 14 animals; median distances were 2.67 and $11.43\,\mu\mathrm{m}$, respectively.\\
\textbf{g,} Spatial footprint correlations for the same directed candidates in \textbf{f}. Median correlations were 0.930 for centroid closest candidates and 0.333 for other candidates.\\
\textbf{h,} Joint distribution of centroid distance and spatial footprint correlation after alignment for 21,412 CellReg registered pairs. Color denotes logarithmically scaled raw bin count and the marginal histograms show linear raw counts. Median centroid distance was $1.74\,\mu\mathrm{m}$ and median footprint correlation was 0.959.\\
\textbf{i,} Distribution of CellReg registration scores for the same registered pair population. The median was 0.997 and 74.7\% were at least 0.95.
}
\labeltoimage{fig:ed04}


\clearpage

\nolinenumbers
\begin{center}
\saveimageanchor{fig:ed05}
\includegraphics[
    width=\textwidth,
    height=0.95\textheight,
    keepaspectratio
]{figures/ed05_version_record.pdf}
\end{center}

\clearpage
\captionof{extendeddatafigure}{
\textbf{Examples and controls for long gap neural stability.}\\
\textbf{a,} Illustrative place map pairs from neurons matched across 3 month gap session pairs. Each map was independently normalized to its 99th percentile rate for display; color denotes normalized rate, and the Pearson correlation between paired maps is shown above each example. Red outlines identify the four examples shown in ~\ref{fig:figure2}b.\\
\textbf{b,} Long gap place map correlations across four day combinations. The labels d1--d1, d1--d2, d2--d1 and d2--d2 denote Day $i$ versus Day $i+90$, Day $i$ versus Day $i+91$, Day $i+1$ versus Day $i+90$, and Day $i+1$ versus Day $i+91$, respectively. Points show animal means, bars show means across animals and error bars show 95\% confidence intervals (36 session pairs from 14 animals). Place map correlation means were 0.12, 0.11, 0.11 and 0.12 across the four day combinations, respectively. No significant differences were detected between d1--d1 and the other comparisons (BH corrected $p$ values were 0.20, 0.58 and 0.88, respectively). \\
\textbf{c,} Same as \textbf{b} but for population vector correlation. Population vector correlations were highest for d1--d1 comparisons (means, 0.22, 0.19, 0.18 and 0.18). d1--d1 exceeded d1--d2, d2--d1 and d2--d2 (BH corrected $p$ values were all $2.4\times10^{-4}$).\\
\textbf{d,} Same as \textbf{b} but for ensemble rate correlation. Ensemble rate correlations were also highest for d1--d1 comparisons (means, 0.53, 0.46, 0.44 and 0.44). d1--d1 exceeded d1--d2, d2--d1 and d2--d2 (BH corrected $p$ values were $3.7\times10^{-4}$, $1.2\times10^{-3}$ and $1.2\times10^{-3}$, respectively).\\
\textbf{e,} Ensemble rate correlations outside the maze and during maze behavior with equal samples for d1--d1 long gap pairs. Gray lines connect the same session pair observations; black points show means across pairs and error bars show 95\% confidence intervals from regressions with standard errors clustered by animal.  Correlation was 0.36 outside the maze (95\% CI, 0.29 to 0.44) and 0.50 inside the maze (0.42 to 0.58). The paired difference was 0.14 (36 pairs from 14 animals; BH corrected $p=1.6\times10^{-5}$).

}
\labeltoimage{fig:ed05}


\clearpage

\nolinenumbers
\begin{center}
\saveimageanchor{fig:ed06}
\includegraphics[
    width=\textwidth,
    height=0.95\textheight,
    keepaspectratio
]{figures/ed06_version_record.pdf}
\end{center}

\clearpage
\captionof{extendeddatafigure}{
\textbf{Convergence within recording cycles and local stability within sessions.}\\
\textbf{a,} Schematic of the fixed reference analysis. Successive windows from the Day $i$ session were compared with the full Day $i+1$ representation from the same recording cycle.\\
\textbf{b,} Schematic of the sliding pivot analysis. Two abutting windows before and after a pivot were compared as the pivot moved through Day $i$ or Day $i+1$.\\
\textbf{c,} Place map correlation between Day $i$ windows and the full Day $i+1$ representation used as the template, using windows of 10 min, 20 min or 3,000 tile transitions.\\ %Slopes were 0.10 $r$ per hour for 10 min windows and 0.09 $r$ per hour for 20 min windows. The slope for windows of 3,000 transitions was 0.02 $r$ per 1,000 transitions. Corresponding BH corrected $p$ values were $5.6\times10^{-13}$, $9.6\times10^{-13}$ and $1.1\times10^{-11}$.
\textbf{d,} Place map correlation between the 30 min windows before and after a sliding pivot. Whole-curve variability was lower on Day $\mathrm{i}+1$ than on Day $\mathrm{i}$ (within-curve s.d., 0.024 versus 0.049 $r$; BH corrected $p=1.4\times10^{-7}$). Separate linear tests on Day $\mathrm{i}$ show a 0.12 $r$ per hour increase (corrected $p=3.6\times10^{-9}$), while Day $\mathrm{i}+1$ shows no significant increase (corrected $p=0.57$). \\
\textbf{e,} Population vector correlation between Day $\mathrm{i}$ windows and the full Day $\mathrm{i}+1$ representation used as the template. \\ %Slopes were 0.104 $r$ per hour for 10 min windows and 0.0941 $r$ per hour for 20 min windows. The slope for windows of 3,000 transitions was 0.0171 $r$ per 1,000 transitions. Corresponding BH corrected $p$ values were $2.5\times10^{-12}$, $5.3\times10^{-12}$ and $3.9\times10^{-9}$.
\textbf{f,} Population vector correlation across the 30 min sliding pivot. Whole-curve variability was lower on Day $\mathrm{i}+1$ than on Day $\mathrm{i}$ (within-curve s.d., 0.027 versus 0.043 $r$; BH corrected $p=3.3\times10^{-5}$). Separate linear tests on Day $\mathrm{i}$ show a 0.09 $r$ per hour increase (corrected $p=3.1\times10^{-8}$) and Day $\mathrm{i}+1$ shows a marginal increase of 0.02 $r$ per hour (corrected $p=0.055$).\\
\textbf{g,} Ensemble rate correlation between Day $\mathrm{i}$ windows and the full Day $\mathrm{i}+1$ representation used as the template. \\%Slopes were 0.0597 $r$ per hour for 10 min windows and 0.0635 $r$ per hour for 20 min windows. The slope for windows of 3,000 transitions was 0.0135 $r$ per 1,000 transitions. Corresponding BH corrected $p$ values were $2.8\times10^{-6}$, $2.1\times10^{-6}$ and $2.4\times10^{-5}$.
\textbf{h,} Ensemble rate correlation across the 30 min sliding pivot. Whole-curve variability did not differ between Day $\mathrm{i}+1$ and Day $\mathrm{i}$ (within-curve s.d., 0.037 versus 0.036 $r$; BH corrected $p=0.88$). Separate linear tests on Day $\mathrm{i}$ show a 0.04 $r$ per hour increase (corrected $p=0.0021$), while Day $\mathrm{i}+1$ shows no significant increase (corrected $p=0.32$). Across \textbf{c--h}, color opacity indicates cumulative experience. Curves show means across animals $\pm$ s.e.m. within each recording cycle. Fixed reference and sliding pivot analyses both used 59 recording cycles from 22 animals. 



}
\labeltoimage{fig:ed06}


\clearpage

\nolinenumbers
\begin{center}
\saveimageanchor{fig:ed07}
\includegraphics[
    width=\textwidth,
    height=0.95\textheight,
    keepaspectratio
]{figures/ed07_version_record.pdf}
\end{center}

%\clearpage
\captionof{extendeddatafigure}{
\textbf{Age associations with map stability.}\\
\textbf{a--c,} Place map, population vector and ensemble rate correlations between the first 30 min of Day $i$ and the full Day $i+1$ representation, plotted against age on Day $i$. The analysis included 59 observations from 22 animals, each representing one animal in one recording cycle. Solid circles show observations, colors indicate cycle stage, and lines show a common age slope with a separate intercept for each cycle stage. Per 30 d of age, place map correlation decreased by 0.012 (95\% CI, $-0.021$ to $-0.0031$; BH corrected $p=0.033$). Population vector correlation decreased by 0.011 ($-0.022$ to $0.000$; corrected $p=0.069$). Ensemble rate correlation decreased by 0.007 ($-0.020$ to 0.008; corrected $p=0.37$). Only the place map correlation decreased significantly.\\
\textbf{d--f,} Long gap place map, population vector and ensemble rate correlations plotted against age in the destination session. The analysis included 36 pairs from 14 animals. Solid circles are session pairs, colors denote destination cycle, and lines show the age association adjusted for transition. Per 30 d, the slopes were $-0.002$ for place maps (95\% CI, $-0.020$ to 0.017; BH corrected $p=0.88$) and $+0.010$ for PVs ($-0.011$ to 0.030; corrected $p=0.63$). The ensemble rate slope was $+0.029$ (0.006 to 0.052; corrected $p=0.066$). No association was significant after correction.

}
\labeltoimage{fig:ed07}


\clearpage

\nolinenumbers
\begin{center}
\saveimageanchor{fig:ed08}
\includegraphics[
    width=\textwidth,
    height=0.95\textheight,
    keepaspectratio
]{figures/ed08_version_record.pdf}
\end{center}

\clearpage
\captionof{extendeddatafigure}{
\textbf{Reward code stability across sampling and route controls.}\\
\textbf{a,} Cross session decoding error map in which each of 14 eligible animals contributed equally. The map was derived from 30 Day $i$ session pairs whose source session had rewarded experience. The map includes the 104 spatial bins supported by every animal. The water drop icon marks the water port; the orange outline marks the reward adjacent bin used in \textbf{b}, and gray outlines mark the matched control bins.\\
\textbf{b,} Reward decoding error compared with the mean error at matched control bins after equalizing either recording frames or visits in both source and target sessions. With Day $i$ sources, differences between control and reward bins were 2.5 bins with equal frames (95\% CI, 1.4 to 3.6; 30 pairs from 14 animals; BH corrected $p=7.3\times10^{-4}$). With equal visits, the difference was 1.9 bins (0.9 to 2.9; corrected $p=1.1\times10^{-3}$). With Day $i+1$ sources, the corresponding differences were 3.0 bins (2.0 to 4.1; 36 pairs from 14 animals; corrected $p=1.2\times10^{-4}$) and 2.7 bins (1.7 to 3.6; corrected $p=1.2\times10^{-4}$). Small points and connecting lines are animal means; large circles and squares with black error bars show means across animals and 95\% confidence intervals.\\
\textbf{c,} Reward and control location population vector correlations across a 3 month gap. Filled circles with solid connectors show same location correlations across reward and control locations; the difference was 0.12 on Day $i$ (95\% CI, 0.08 to 0.16; 30 pairs from 14 animals; BH corrected $p=8.2\times10^{-5}$). On Day $i+1$, the difference was 0.10 (0.06 to 0.14; 36 pairs from 14 animals; corrected $p=1.6\times10^{-4}$). Open circles with dashed connectors show the cross location correlations. There is no evidence of a significant difference on both Day $i$ and Day $i+1$. Confidence intervals and two sided tests used standard errors clustered by animal. Large symbols and error bars show means and intervals.\\
\textbf{d,} Reward end advantage of the population vector correlation in individual animals: the correlation at the water port minus the mean correlation elsewhere on the route, for home to water (orange diamonds) and water to home runs (blue squares), for the sessions in Fig.~\ref{fig:figure4}e (filled) and for the second day of each visit, Day~$i+1$ versus Day~$i+91$ (open). For the Fig.~\ref{fig:figure4}e sessions, the advantage was positive in all 13 animals on home to water runs and in 11 of 14 animals on water to home runs. On the second day it was 0.11 on home to water runs (95\% CI, 0.03 to 0.18; BH corrected $p=0.0037$; positive in 10 of 13 animals; 29 session pairs) and 0.09 on water to home runs (0.04 to 0.14; corrected $p=0.0037$; positive in 11 of 14 animals; 28 pairs; two sided exact sign-flip tests with BH correction across the two directions). Solid and dotted vertical lines mark the means across animals for the Fig.~\ref{fig:figure4}e sessions and for the second day, and the gray dashed line marks zero.\\
\textbf{e,} Reliability-adjusted population vector correlation: the across-gap correlation between half-session maps divided by the within-day correlation between maps from odd and even runs, at the water port and elsewhere on the route. The reward end advantage remained after accounting for within-day reliability: the reliability-adjusted correlation was 0.28 at the water port and 0.14 elsewhere on home to water runs (difference, 0.14; 95\% CI, 0.08 to 0.20; BH corrected $p=9.8\times10^{-4}$), and 0.28 and 0.16 on water to home runs (difference, 0.13; 0.01 to 0.24; corrected $p=0.035$). Both means and the paired test used the same 13 animals for home to water and 13 for water to home, with defined ratios at both locations.\\
\textbf{f,} Equal frame control. In each of 100 draws, every route bin was sampled to the session's smallest bin occupancy, so that every position contributed the same number of frames. The correlation at the water port exceeded the mean elsewhere by 0.13 on home to water runs (95\% CI, 0.09 to 0.16; BH corrected $p=6.5\times10^{-4}$) and by 0.09 on water to home runs (0.04 to 0.14; corrected $p=0.0057$).\\
\textbf{g,} Route coordinates for runs between the water port and the home, node 2 or node 3. \\
\textbf{h,} Population vector correlation at the water port and elsewhere on runs from nodes 2 and 3 to the water port and on the reverse runs from the water port (30 session pairs from 14 animals for node 2 and 29 pairs from 14 animals for node 3). The water port exceeded the rest of the route by 0.13 on runs from node 2 to the water port and by 0.11 on the reverse runs (BH corrected $p=4.9\times10^{-4}$ and $0.0014$), and by 0.12 and 0.09 on the same two runs for node 3 (corrected $p=4.9\times10^{-4}$ and $0.0018$).\\
\textbf{i,} Reward end advantage of the population vector correlation on rewarded and timeout home to water runs (25 session pairs from 12 animals). The advantage was present on both (rewarded, 0.18; 95\% CI, 0.09 to 0.26; timeout, 0.18; 0.13 to 0.23; BH corrected $p=4.9\times10^{-4}$ and $4.9\times10^{-4}$), and the rewarded minus timeout difference was 0.002 (95\% CI, $-0.083$ to 0.086; $p=0.98$).\\
\textbf{j,} Population vector correlation between Day~$i$ and Day~$i+90$ at each route position for the session pairs analyzed in both directions (27 pairs from 13 animals; the same neurons for every curve): home to water runs on both days (orange), water to home runs on both days (blue), and approach on one day against departure on the other (black; shading, 95\% CI), with the neuron ID shuffle of the approach versus departure comparison (dotted). At the water port, approach and departure patterns correlated at 0.17, against 0.02 elsewhere on the route (difference, 0.15; 95\% CI, 0.06 to 0.23; two sided exact sign-flip test, $p=0.0032$). In \textbf{e}, \textbf{f}, \textbf{h} and \textbf{i}, small points and connecting lines are animal values, and large symbols and error bars show means across animals and 95\% CIs. Mean neuron ID shuffle correlations were within 0.003 of zero at the water port and elsewhere in \textbf{d}, \textbf{f}, \textbf{h}, \textbf{i} and \textbf{j}. All tests were two sided exact sign-flip tests across animals, with BH correction across the eight comparisons in \textbf{e}, \textbf{f} and \textbf{h} and across the two run types in \textbf{i}. 

}
\labeltoimage{fig:ed08}


\clearpage

\nolinenumbers
\begin{center}
\saveimageanchor{fig:ed09}
\includegraphics[
    width=\textwidth,
    height=0.95\textheight,
    keepaspectratio
]{figures/ed09_version_record.pdf}
\end{center}

\clearpage
\captionof{extendeddatafigure}{
\textbf{Additional analyses showing persistence of structure code.}\\
\textbf{a,} Preferred level transitions among the 12,541 matched neuron observations that changed level, from 16,901 observations across 36 session pairs and 14 animals. For each Day $i$ level, the heat map shows the observed probability of each Day $i+90$ level minus the mean probability from 10,000 neuron ID permutations within each pair. Observed and shuffled probabilities were separately normalized across off diagonal Day $i+90$ levels.\\
\textbf{b,} Preferred level shift distance among the same changed level neurons. Adjacent level shifts exceeded the permutation expectation by 0.069 (95\% CI, 0.051 to 0.087; BH corrected $p=1.6\times10^{-6}$). Mean shifts were 0.19 levels shorter than expected (0.14 to 0.25; corrected $p=2.8\times10^{-6}$). The solid curve shows the observed distribution; the dashed curve and error bars show the permutation null mean $\pm$ s.d.\\
\textbf{c,} Long gap place map correlation compared with a shuffle that preserved tree depth but permuted fine position within each level. Across 36 session pairs from 14 animals, the Data mean was 0.12 and the shuffle mean was 0.06. Their difference was 0.06 (95\% CI, 0.01 to 0.10; one sided $p=0.0074$ from an intercept only regression with standard errors clustered by animal). Left, thin lines connect each pair's gray shuffle and red data values; large symbols show means across pairs. Right, the gray histogram shows 1,000 cohort mean correlations under the within-level shuffle, and the red line marks the data mean.\\
\textbf{d,} Schematic of the geometric eight bin levels and the horizontal and vertical eight bin spatial partitions with the same resolution.\\
\textbf{e,} Average tree depth seven bin and geometric eight bin level map correlations (36 pairs from 14 animals). Means were 0.25 and 0.24, across session pairs.\\
\textbf{f,} Average correlations for horizontal, geometric and vertical eight bin maps in the 36 session pairs from 14 animals. Means were 0.09, 0.24 and 0.08, respectively across session pairs, with BH corrected $p$ values of $1.2\times10^{-5}$ and $3.7\times10^{-7}$. Small connected circles show session pairs and large circles show cohort means.\\
\textbf{g,} Density of joint long gap place map and level map correlations for 16,901 matched neuron observations from 36 session pairs and 14 animals.\\
\textbf{h,} Level decoding accuracy minus neuron ID shuffle accuracy after excluding neurons whose place map correlation exceeded the displayed threshold. At the analysis threshold of 0.3, the difference was 0.07 (95\% CI, 0.06 to 0.09; 36 pairs from 14 animals; one sided BH corrected $p=1.7\times10^{-7}$). Points show observed minus shuffle estimates; shading shows 95\% confidence intervals from regressions with standard errors clustered by animal. The dashed horizontal line marks zero and the vertical red line marks the analysis threshold of 0.3.\\
\textbf{i,} Level map correlation across intervals of 1, 90, 180, 270 and 360 d. Values are means across session pairs with 95\% confidence intervals. The dashed curve shows neuron ID shuffle means. The 1 day point compares Day $i$ with Day $i+1$ within the same recording cycle (47 Day $i$ to Day $i+1$ session pairs from 14 animals). The longer intervals compare Day $i$ sessions from successive recording cycles (36, 22, 11 and 4 session pairs from 14, 11, 7 and 4 animals, respectively).\\
\textbf{j,} Level map stability across consecutive long gaps. The analysis comprised 9,661 matched neuron observations from 22 three session chains and 11 animals. A linear regression with standard errors clustered by animal gave a slope of 0.18 (BH corrected $p=3.6\times10^{-4}$). Previous gap correlations were divided into equal width bins across the x axis. The solid curve connects the mean next gap correlation at each bin center. Faint points show the underlying matched neuron observations.\\
\textbf{k,} Tree depth level map correlation between moving 30 min Day $i$ windows and the full Day $i+1$ representation used as the template. Shaded bands show the smoothed correlation s.e.m. Color opacity indicates cumulative experience. The analysis comprised 59 recording cycles from 22 animals, with 22, 14, 12, 7 and 4 animals contributing at naive, 3, 6, 9 and 12 month experience, respectively.\\
\textbf{l,} Long gap level map correlation between Day $i$ sessions from successive recording cycles (37 session pairs from 14 animals, including one short interval pair). Correlation increased with experience (Pearson $r=0.53$; BH corrected $p=0.034$). Points denote session pairs. The line and shading show the fitted relation and its 95\% confidence interval; inference used standard errors clustered by animal.

}
\labeltoimage{fig:ed09}


\clearpage

\nolinenumbers
\begin{center}
\saveimageanchor{fig:ed10}
\includegraphics[
    width=\textwidth,
    height=0.95\textheight,
    keepaspectratio
]{figures/ed10_version_record.pdf}
\end{center}

\clearpage
\captionof{extendeddatafigure}{
\textbf{Additional neural vs behavior correlation analysis.}\\
\textbf{a,} Pearson correlations between six neural readouts from session pairs with a 3 month gap and three behaviors measured on Day $i$ of the later session (30 session pairs from 14 animals). Reward approach speed advantage is the log ratio of median center of mass speed when approaching rewarded end node versus end nodes at the same depth that were not rewarded. All neural readouts used the same number of neurons from all session pairs. Both decoder rows used exactly 2,000 frames per pair. Stars represent p values with BH correction.\\
\textbf{b,} Commonality partition of variance jointly associated with reward decoding error and level map correlation. Across the same 30 session pairs from 14 animals, the two neural measures contributed differently to the two behavior measures (two sided sign-flip test, BH corrected $p=0.013$). Reward decoding was more strongly associated with earlier long bouts than level map correlation was, but this difference was not significant (difference in r, 0.12; corrected $p=0.33$). Conversely, level map correlation was more strongly associated with exploration efficiency than reward decoding was (difference in r, 0.28; corrected $p=0.013$). \\
\textbf{c,} Level map correlation plotted against four turning policies measured from exploration mode during the first 30 min of time in maze. This analysis comprised 36 session pairs from the same 14 animals. The depth weighted $\mathrm{P}_{\mathrm{BS}}$ weights each event by $2^{6-d}$, where $d$ is decision node depth. The depth weighted $\mathrm{P}_{\mathrm{BS}}$ correlation was supported by a two sided regression with standard errors clustered by animal (BH corrected $p=1.4\times10^{-5}$). 

}
\labeltoimage{fig:ed10}


\clearpage
\includepdf[pages=-,angle=90,fitpaper=false,pagecommand={\thispagestyle{empty}\hypertarget{supp-table-1}{}}]{figures/Supplementary_Table_S1_Mouse_Cohorts_by_Analysis.pdf}
\includepdf[pages=-,angle=90,fitpaper=false,pagecommand={\thispagestyle{empty}\hypertarget{supp-table-2}{}}]{figures/Supplementary_Table_2_Long_Gap_Intervals.pdf}

\end{document}
